\frfilename{mt284.tex} \versiondate{30.8.13} \copyrightdate{1994}

\def\rdtf{fungsi uji yang menurun cepat} \def\Ft{transformasi Fourier}

\def\chaptername{Analisis Fourier} \def\sectionname{Transformasi Fourier II}

\newsection{284}

Paradoks dasar transformasi Fourier adalah bahwa meskipun untuk fungsi-fungsi
tertentu (lihat 283J-283K) kita mempunyai $(\varhatf)\varspcheck=f$,
fungsi terintegralkan `biasa' $f$ (misalnya fungsi indikator interval
nontrivial) menghasilkan transformasi Fourier $\varhatf$ yang tidak dapat
diintegralkan. Karena tidak tersedia definisi langsung untuk
$\varhatf\varcheck{\phantom{f}}$, kita harus menentukan dalam arti apa rumus
$f=\varhatf\varcheck{\phantom{f}}$ mungkin benar. Penyelesaian yang kini
tampak paling alami ialah menyatakan transformasi Fourier sebagai operasi pada
{\it distribusi}, bukan pada {\it fungsi}. Saya tidak akan mencoba menguraikan
teori ini secara lengkap (hampir setiap buku tentang `Distribusi' akan
membahas pokok ini lebih baik daripada yang dapat saya lakukan di sini), tetapi
akan menyampaikan gagasan-gagasan dasarnya sejauh relevan dengan pertanyaan di
sini, dengan bahasa yang memudahkan peralihan ke pembahasan yang lebih lengkap.
Pada saat yang sama, metode ini memudahkan pembuktian versi-versi kuat teorema
`klasik' mengenai transformasi Fourier.

\leader{284A}{Fungsi uji: Definisi} Sepanjang bagian ini, {\bf fungsi uji yang
menurun cepat} atau {\bf fungsi Schwartz} berarti fungsi $h:\Bbb R\to\Bbb C$
sedemikian rupa sehingga $h$ {\bf mulus}, yaitu semua turunannya hingga orde
berapa pun ada di setiap titik, dan selain itu

\Centerline{$\sup_{x\in\Bbb R}|x|^k|h^{(m)}(x)|<\infty$}

\noindent untuk semua $k$, $m\in\Bbb N$, dengan $h^{(m)}$ menyatakan turunan
ke-$m$ dari $h$.

\leader{284B}{}\cmmnt{ Fakta dasar berikut akan berguna.

\medskip

\noindent}{\bf Lema} (a) Jika $g$ dan $h$ adalah fungsi-fungsi uji yang
menurun cepat, maka juga $g+h$ dan $ch$, untuk setiap $c\in\Bbb C$.

(b) Jika $h$ adalah fungsi uji yang menurun cepat dan $y\in\Bbb R$, maka
$x\mapsto h(y-x)$ adalah fungsi uji yang menurun cepat.

(c) Jika $h$ adalah fungsi uji yang menurun cepat, maka $h$ dan $h^2$ adalah
terintegralkan.

(d) Jika $h$ adalah \rdtf, maka turunannya $h'$ juga demikian.


(e) Jika $h$ adalah \rdtf, maka fungsi $x\mapsto xh(x)$ juga demikian.

(f) Untuk setiap $\epsilon>0$, fungsi $x\mapsto e^{-\epsilon x^2}$ adalah
fungsi uji yang menurun cepat.

\proof{{\bf (a)} Jelas.

\medskip

{\bf (b)} Tulis $g(x)=h(y-x)$ untuk $x\in\Bbb R$. Kemudian
$g^{(m)}(x)=(-1)^mh^{(m)}(y-x)$ untuk setiap $m$, jadi $g$ halus. Untuk setiap
$k\in\Bbb N$,

\Centerline{$|x|^k\le 2^k(|y|^k+|y-x|^k)$}

\noindent untuk setiap $x$, jadi

$$\eqalign{\sup_{x\in\Bbb R}|x|^k|g^{(m)}(x)|
&=\sup_{x\in\Bbb R}|x|^k|h^{(m)}(y-x)|\cr
&\le 2^k|y|^k\sup_{x\in\Bbb R}|h^{(m)}(y-x)|
   +2^k\sup_{x\in\Bbb R}|y-x|^k|h^{(m)}(y-x)|\cr
&=2^k|y|^k\sup_{x\in\Bbb R}|h^{(m)}(x)|
   +2^k\sup_{x\in\Bbb R}|x|^k|h^{(m)}(x)|
<\infty.\cr}$$

\medskip

{\bf (c)} Karena

\Centerline{$M=\sup_{x\in\Bbb R}(|h(x)|+x^2|h(x)|)$}

\noindent adalah hingga, kita memiliki

\Centerline{$\int|h|\le\int\Bover{M}{1+x^2}dx<\infty$.}

\noindent Tentu saja sekarang kita mempunyai $|h^2|\le M|h|$, sehingga $h^2$
juga terintegralkan.

\medskip

{\bf (d)} Ini langsung dari definisi, karena setiap turunan dari $h'$ adalah
turunan dari $h$.

\medskip

{\bf (e)} Dengan menetapkan $g(x)=xh(x)$, kita memperoleh
$g^{(m)}(x)=xh^{(m)}(x)+mh^{(m-1)}(x)$ untuk $m\ge 1$, sehingga

\Centerline{$\sup_{x\in\Bbb R}|x^kg^{(m)}(x)| \le\sup_{x\in\Bbb
R}|x^{k+1}h^{(m)}(x)| +m\sup_{x\in\Bbb R}|x^kh^{(m-1)}(x)|$}

\noindent adalah hingga, untuk semua $k\in\Bbb N$, $m\ge 1$.

\medskip

{\bf (f)} Jika $h(x)=e^{-\epsilon x^2}$, maka untuk setiap $m\in\Bbb N$ kita
mempunyai $h^{(m)}(x)=p_m(x)h(x)$, dengan $p_0(x)=1$ dan
$p_{m+1}(x)=p'_m(x)-2\epsilon xp_m(x)$, sehingga $p_m$ adalah polinom. Karena
$e^{\epsilon x^2}\ge\epsilon^{k+1}x^{2k+2}/(k+1)!$ untuk semua $x$ dan
$k\ge 0$,

\Centerline{$\lim_{|x|\to\infty}|x|^kh(x) =\lim_{x\to\infty}x^k/e^{\epsilon
x^2}=0$}

\noindent untuk setiap $k$, dan $\lim_{|x|\to\infty}p(x)h(x)=0$ untuk setiap
polinomial $p$; karenanya

\Centerline{$\lim_{|x|\to\infty}x^kh^{(m)}(x)
=\lim_{|x|\to\infty}x^kp_m(x)h(x)=0$}

\noindent untuk semua $k$, $m$. Jadi, $h$ adalah fungsi uji yang menurun cepat. }%end of proof of 284B


\leader{284C}{Proposisi} Misalkan $h:\Bbb R\to\Bbb C$ adalah fungsi uji yang
menurun cepat. Maka $\varhat{h}:\Bbb R\to\Bbb C$ dan $\varcheck{h}:\Bbb
R\to\Bbb C$ juga merupakan fungsi uji yang menurun cepat, dan
$\varhat{h}\varcheck{\phantom{h}} = \varcheck{h}\varhat{\phantom{h}} = h$.

\proof{{\bf (a) } Ambil $k$, $m\in\Bbb N$. Maka $\sup_{x\in\Bbb
R}(|x|^m+|x|^{m+2})|h^{(k)}(x)|<\infty$ dan
$\int_{-\infty}^{\infty}|x^mh^{(k)}(x)|dx
\ifdim\pagewidth=390pt\penalty-100\fi <\infty$. Oleh karena itu, menurut
283Ch-283Ci, $y\mapsto
i^{k+m}y^k\varhat{h}\vthsp^{(m)}(y)$ adalah transformasi Fourier dari
$x\mapsto x^mh^{(k)}(x)$. Menurut 283Cg, karena itu
$\lim_{|y|\to\infty}y^k\varhat{h}\vthsp^{(m)}(y)=0$, sehingga
(karena $\varhat{h}\vthsp^{(m)}$ adalah kontinu) $\sup_{y\in\Bbb
R}|y^k\varhat{h}\vthsp^{(m)}(y)|$ adalah hingga. Karena $k$ dan $m$ sebarang,
$\varhat{h}$ adalah fungsi uji yang menurun cepat.

\medskip

{\bf (b)} Karena $\varcheck{h}(y)=\varhat{h}(-y)$ untuk setiap $y$, langsung
terlihat bahwa $\varcheck{h}$ adalah fungsi uji yang menurun cepat.

\medskip

{\bf (c)} Menurut 283J, (a) dan (b) memberikan
$\varhat{h}\varcheck{\phantom{h}}=\varcheck{h}\varhat{\phantom{h}}=h$. }%end of proof of 284C

\leader{284D}{Definisi} Saya akan menggunakan frasa {\bf fungsi tempered} pada
$\Bbb R$ untuk menyatakan fungsi terukur bernilai kompleks $f$, yang
terdefinisi hampir di mana-mana dalam $\Bbb R$ dan memenuhi

\Centerline{$\int_{-\infty}^{\infty} \Bover1{1+|x|^k}|f(x)|dx<\infty$}

\noindent untuk beberapa $k\in\Bbb N$.

\leader{284E}{}\cmmnt{ Seperti dalam 284B, berikut beberapa fakta dasar.

\medskip

\noindent}{\bf Lema} (a) Jika $f$ dan $g$ adalah fungsi-fungsi tempered, maka
juga $|f|$, $f+g$ dan $cf$, untuk setiap $c\in\Bbb C$.

(b) Jika $f$ adalah fungsi tempered, maka fungsi tersebut terintegralkan pada
setiap interval terbatas.

(c) Jika $f$ adalah fungsi tempered dan $x\in\Bbb R$, maka $t\mapsto f(x+t)$
dan $t\mapsto f(x-t)$ keduanya merupakan fungsi tempered.

\proof{{\bf (a)} Ini bersifat elementer; jika

\Centerline{$\int_{-\infty}^{\infty} \Bover1{1+|x|^j}|f(x)|dx<\infty$,
\quad$\int_{-\infty}^{\infty}\Bover1{1+|x|^k}|g(x)|dx<\infty$,}

\noindent maka

\Centerline{$\int_{-\infty}^{\infty}\Bover1{1+|x|^{j+k}} |(f+g)(x)|dx<\infty$}

\noindent karena

\Centerline{$1+|x|^{j+k}\ge\max(1,|x|^{j+k})\ge\max(1,|x|^j,|x|^k)
\ge\Bover12\max(1+|x|^j,1+|x|^k)$}

\noindent untuk semua $x$.

\medskip

{\bf (b)} Jika

\Centerline{$\int_{-\infty}^{\infty}\Bover1{1+|x|^k}|f(x)|dx =M<\infty$,}

\noindent maka untuk $a\le b$

\Centerline{$\int_a^b|f|\le M(1+|a|^k+|b|^k)(b-a)<\infty$.}

\medskip

{\bf (c)} Gagasannya sama seperti dalam 284Bb. Misalkan $k\in\Bbb N$
sedemikian sehingga

\Centerline{$\int_{-\infty}^{\infty}\Bover1{1+|t|^k}|f(t)|dt =M<\infty$,}

\noindent maka kita memiliki

\Centerline{$1+|x+t|^k \le 2^k(1+|x|^k)(1+|t|^k)$}

\noindent sehingga

\Centerline{$\Bover1{1+|t|^k} \le 2^k(1+|x|^k)\Bover1{1+|x+t|^k}$}

\noindent untuk setiap $t$, dan

\Centerline{$\int_{-\infty}^{\infty}\Bover{|f(x+t)|}{1+|t|^k}dt \le
2^k(1+|x|^k)\int_{-\infty}^{\infty} \Bover{|f(x+t)|}{1+|x+t|^k}dt \le
2^k(1+|x|^k)M<\infty$.}

\noindent Sama halnya,

\Centerline{$\int_{-\infty}^{\infty}\Bover{|f(x-t)|}{1+|t|^k}dt
\le 2^k(1+|x|^k)M<\infty$.}
}%end of proof of 284E

\leader{284F}{}\cmmnt{ Dengan menghubungkan kedua konsep tersebut, kita
memperoleh hasil berikut.

\medskip

\noindent}{\bf Lema} Misalkan $f$ adalah fungsi tempered pada $\Bbb R$ dan
$h$ adalah fungsi uji yang menurun cepat. Maka $f\times h$ terintegralkan.

\proof{ Tentu saja $f\times h$ terukur. Pilih $k\in\Bbb N$ sedemikian sehingga
$\int_{-\infty}^{\infty} \bover1{1+|x|^k}|f(x)|dx<\infty$. Ada $M$ sedemikian
sehingga $(1+|x|^k)|h(x)|\le M$ untuk setiap $x\in\Bbb R$, sehingga

\Centerline{$\int_{-\infty}^{\infty}|f\times h|
\le M\int_{-\infty}^{\infty}\Bover1{1+|x|^k}|f(x)|dx<\infty$.}
}%end of proof of 284F

\leader{284G}{Lema} Misalkan $f_1$ dan $f_2$ adalah fungsi-fungsi tempered dan
$\int f_1\times h=\int f_2\times h$ untuk setiap fungsi uji yang menurun cepat
$h$. Maka $f_1\eae f_2$.

\proof{{\bf (a)} Tetapkan $g=f_1-f_2$; maka $\int g\times h=0$ untuk setiap
fungsi uji yang menurun cepat $h$. Tentu saja $g$ adalah fungsi tempered, jadi
terintegralkan pada setiap interval terbatas. Menurut 222D, cukup ditunjukkan
bahwa $\int_a^bg=0$ setiap kali $a<b$; dengan demikian $g=0$ a.e.\ pada setiap
interval terbatas dan $f_1\eae f_2$.

\medskip

{\bf (b)} Perhatikan fungsi $\tilde\phi(x)=e^{-1/x}$ untuk $x>0$. Fungsi
$\tilde\phi$ dapat diturunkan hingga orde berapa pun di setiap titik dalam
$\ooint{0,\infty}$, $0<\tilde\phi(x)<1$ untuk setiap $x>0$, dan
$\lim_{x\to\infty}\tilde\phi(x)=1$. Selain itu, menulis $\tilde\phi^{(m)}$
untuk turunan ke-$m$ dari $\tilde\phi$,

\Centerline{$\lim_{x\downarrow 0}\tilde\phi^{(m)}(x) =\lim_{x\downarrow
0}\Bover1x\tilde\phi^{(m)}(x)=0$}

\noindent untuk setiap $m\in\Bbb N$. \Prf\ (Bandingkan 284Bf.) Kita mempunyai
$\tilde\phi^{(m)}(x)=p_m(\bover1x)\tilde\phi(x)$, di mana $p_0(t)=1$ dan
$p_{m+1}(t)=t^2(p_m(t)-p_m'(t))$, sehingga $p_m$ adalah polinom untuk setiap
$m\in\Bbb N$. Sekarang, untuk setiap $k\in\Bbb N$,

\Centerline{$0\le\limsup_{t\to\infty}t^ke^{-t}
\le\lim_{t\to\infty}\Bover{(k+1)!t^k}{t^{k+1}}=0$,}

\noindent jadi

\Centerline{$\lim_{x\downarrow 0}\tilde\phi^{(m)}(x)
=\lim_{t\to\infty}p_m(t)e^{-t}=0$,}

\Centerline{$\lim_{x\downarrow 0}\Bover1x\tilde\phi^{(m)}(x)
=\lim_{t\to\infty}tp_m(t)e^{-t}=0$. \Qed}

\medskip

{\bf (c)} Akibatnya, jika kita menetapkan $\phi(x)=0$ untuk $x\le 0$ dan
$e^{-1/x}$ untuk $x>0$, maka $\phi$ mulus, dengan turunan ke-$m$

\Centerline{$\phi^{(m)}(x)=0$ untuk $x\le 0$,
\quad$\phi^{(m)}(x)=\tilde\phi^{(m)}(x)$ untuk $x>0$.}

\noindent (Buktinya merupakan induksi sederhana pada $m$.) Juga
$0\le\phi(x)\le 1$
untuk setiap $x\in\Bbb R$, dan $\lim_{x\to\infty}\phi(x)=1$.

\medskip

{\bf (d)} Sekarang ambil $a<b$, dan untuk $n\in\Bbb N$ tetapkan

\Centerline{$\phi_n(x)=\phi(n(x-a))\phi(n(b-x))$.}

\noindent Maka $\phi_n$ akan halus dan $\phi_n(x)=0$ jika
$x\notin\ooint{a,b}$, jadi pasti $\phi_n$ adalah fungsi uji yang menurun
cepat, dan

\Centerline{$\int_{-\infty}^{\infty}g\times\phi_n=0$.}

\noindent Selanjutnya, $0\le\phi_n(x)\le 1$ untuk setiap $x$, $n$, dan jika
$a<x<b$, maka $\lim_{n\to\infty}\phi_n(x)=1$. Jadi

\Centerline{$\int_a^bg=\int g\times\chi(\ooint{a,b}) =\int
g\times(\lim_{n\to\infty}\phi_n) =\lim_{n\to\infty}\int g\times\phi_n=0$,}

\noindent menurut teorema konvergensi terdominasi Lebesgue. Karena $a$ dan $b$
sebarang, $g=0$ a.e., sebagaimana diperlukan. }%end of proof of 284G

\leader{284H}{Definisi} Misalkan $f$ dan $g$ adalah fungsi-fungsi
tempered\cmmnt{ dalam arti 284D}. Kita mengatakan bahwa $g$ {\bf mewakili
transformasi Fourier dari} $f$ jika

\Centerline{$\int_{-\infty}^{\infty}g\times h
=\int_{-\infty}^{\infty}f\times\varhat{h}$}

\noindent untuk setiap fungsi uji yang menurun cepat $h$.

\leader{284I}{Catatan (a)}\cmmnt{ Seperti biasa, perubahan sudut pandang dalam
definisi ini memerlukan beberapa pemeriksaan.} Jika $f$ adalah fungsi
terintegralkan bernilai kompleks pada $\Bbb R$ dan $\varhatf$ adalah
transformasi Fouriernya, maka \cmmnt{ $\varhatf$ tentu merupakan fungsi
tempered, sebab kontinu dan terbatas; dan jika $h$ adalah fungsi uji yang
menurun cepat, maka $\int\varhatf\times h=\int f\times\varhat{h}$ menurut
283O. Jadi } $\varhatf$ `mewakili transformasi Fourier dari $f$' \cmmnt{
dalam arti 284H } di atas.

\header{284Ib}{\bf (b)} Perhatikan juga bahwa \cmmnt{ menurut 284G, } jika
$g_1$, $g_2$ adalah dua fungsi tempered yang sama-sama mewakili transformasi
Fourier dari $f$, maka $g_1\eae g_2$\cmmnt{, sebab kita mempunyai

\Centerline{$\int g_1\times h=\int f\times\varhat{h}=\int g_2\times h$}

\noindent untuk setiap fungsi uji yang menurun cepat $h$.}

\spheader 284Ic Jelas bahwa jika $f_1$, $f_2$, $g_1$ dan $g_2$ adalah
fungsi-fungsi tempered dan $g_i$ mewakili transformasi Fourier dari $f_i$
untuk kedua nilai $i$, maka $cg_1+g_2$ mewakili transformasi Fourier dari
$cf_1+f_2$
untuk setiap $c\in\Bbb C$.

\cmmnt{\spheader 284Id Nilai pendekatan tak langsung ini tentu saja terletak
pada kemampuannya memberikan transformasi Fourier, dalam suatu arti, kepada
lebih banyak fungsi. Namun harus segera dicatat bahwa jika $g$ `mewakili
transformasi Fourier dari $f$', maka setiap fungsi yang setara dengan $g$
hampir di mana-mana juga demikian; kita tidak lagi dapat berharap berbicara
tentang `transformasi Fourier' dari $f$ sebagai suatu fungsi tertentu. Kita
dapat mengatakan bahwa transformasi Fourier dari $f$ adalah fungsional $\phi$
pada ruang fungsi uji yang menurun cepat, yang didefinisikan dengan
$\phi(h)=\int f\times\varhat{h}$; atau kita dapat mengatakan bahwa transformasi
Fourier dari $f$ adalah anggota $L_{\Bbb C}^0$, yaitu ruang kelas ekuivalensi
fungsi terukur modulo kesamaan hampir di mana-mana (241J). }%end of comment

\spheader 284Ie\cmmnt{ Sekarang wajar untuk mengatakan bahwa } {\bf $g$
mewakili transformasi Fourier invers dari $f$}\cmmnt{ hanya ketika $f$
mewakili transformasi Fourier dari $g$; yaitu, } ketika $\int f\times h=\int
g\times\varhat{h}$ untuk setiap fungsi uji yang menurun cepat $h$.
\cmmnt{Karena
$\varhat{h}\varcheck{\phantom{h}}=\varcheck{h}\varhat{\phantom{h}}=h$ untuk
setiap $h$ (284C), ini adalah hal yang sama dengan mengatakan bahwa $\int
g\times h=\int f\times\varcheck{h}$ untuk setiap fungsi uji yang menurun cepat
$h$, yang merupakan ekspresi alami lain dari apa yang mungkin berarti bahwa
`$g$ mewakili transformasi Fourier invers dari $f$ }.

\spheader 284If Jika $f$, $g$ adalah fungsi-fungsi tempered dan kita menulis
$\Reverse{g}(x)=g(-x)$ setiap kali ini didefinisikan, maka\cmmnt{
$\Reverse{g}$ juga merupakan fungsi tempered, dan kita selalu mempunyai

\Centerline{$\int\Reverse{g}\times\varhat{h} =\int g(-x)\varhat{h}(x)dx =\int
g(x)\varhat{h}(-x)dx =\int g\times\varcheck{h}$,}

\noindent sehingga}

\qquad $g$ mewakili transformasi Fourier dari $f$

\cmmnt{\qquad\qquad$\iff\,\int g\times h=\int f\times\varhat{h}$ untuk setiap
fungsi uji $h$

\qquad\qquad$\iff\,\int g\times\varcheck{h}=\int
f\times\varcheck{h}\varhat{\phantom{h}}$ untuk setiap $h$

\qquad\qquad$\iff\,\int\Reverse{g}\times\varhat{h}=\int f\times h$ untuk setiap $h$ }%end of comment

\qquad\qquad$\iff\,\Reverse{g}$ mewakili transformasi Fourier invers dari $f$.

\cmmnt{\noindent Dengan menggabungkan ini dengan (d), kita memperoleh

\qquad $g$ mewakili \Ft\ dari $f$

%\qquad\qquad$\iff\,(\Reverse{f})\ssplrarrow=f$
%represents the inverse \Ft\ of $g$

\qquad\qquad$\iff\,\Reverse{\Reverse{f}}=f$ mewakili \Ft\ invers dari $g$

\qquad\qquad$\iff\,\Reverse{f}$ mewakili \Ft\  dari $g$. }%end of comment

\cmmnt{\spheader 284Ig Sekali lagi diperlukan pemeriksaan: jika $f$
terintegralkan dan $\varcheckf$ adalah transformasi Fourier inversnya seperti
yang didefinisikan dalam 283Ab, maka

\Centerline{$\int\varcheckf\times\varhat{h} =\int
f\times\varhat{h}\varcheck{\phantom{h}} =\int f\times h$}

\noindent untuk setiap fungsi uji yang menurun cepat $h$, jadi $\varcheckf$ ` mewakili transformasi Fourier invers dari $f$' dalam arti yang diberikan di sini. }%end of comment

\leader{284J}{Lema} Misalkan $f$ adalah fungsi tempered dan $h$ fungsi uji yang
menurun cepat. Maka $f*h$, yang didefinisikan dengan rumus

\Centerline{$(f*h)(y)=\int_{-\infty}^{\infty}f(t)h(y-t)dt$,}

\noindent terdefinisi di setiap titik.

\proof{ Ambil $y\in\Bbb R$. Menurut 284Bb, $t\mapsto h(y-t)$ adalah fungsi uji
yang menurun cepat, sehingga menurut 284F integral tersebut selalu terdefinisi
dalam $\Bbb C$.
}%end of proof of 284J

\leader{284K}{Proposisi} Misalkan $f$ dan $g$ adalah fungsi-fungsi tempered,
$g$ mewakili transformasi Fourier dari $f$, dan $h$ adalah fungsi uji yang
menurun cepat.

(a) Transformasi Fourier dari fungsi terintegralkan $f\times h$ adalah
$\bover1{\sqrt{2\pi}}g*\varhat{h}$.

(b) Transformasi Fourier dari fungsi kontinu $f*h$ diwakili oleh produk
$\sqrt{2\pi}g\times\varhat{h}$.

\proof{{\bf (a)} Tentu saja $f\times h$ terintegralkan menurut 284F,
sementara $g*\varhat{h}$ terdefinisi di setiap titik menurut 284C dan 284J.

Tetapkan $y\in\Bbb R$. Definisikan $h_1(x)=\varhat{h}(y-x)$ untuk $x\in\Bbb R$;
maka
$h_1$ adalah fungsi uji yang menurun cepat karena $\varhat{h}$ adalah (284Bb).
Sekarang

$$\eqalign{\varhat{h}_1(t)
&=\Bover1{\sqrt{2\pi}}\int_{-\infty}^{\infty}
  e^{-itx}\varhat{h}(y-x)dx
=\Bover1{\sqrt{2\pi}}\int_{-\infty}^{\infty}
  e^{-it(y-x)}\varhat{h}(x)dx\cr
&=\Bover{1}{\sqrt{2\pi}}e^{-ity}
   \int_{-\infty}^{\infty}e^{itx}\varhat{h}(x)dx
=e^{-ity}\varhat{h}\varcheck{\phantom{h}}(t)
=e^{-ity}h(t),\cr}$$

\noindent menurut 284C. Dengan demikian,

$$\eqalignno{(f\times h)\varsphat(y)
&=\Bover1{\sqrt{2\pi}}\int_{-\infty}^{\infty}e^{-ity}f(t)h(t)dt\cr
&=\Bover1{\sqrt{2\pi}}
  \int_{-\infty}^{\infty}f(t)\varhat{h}_1(t)dt
=\Bover1{\sqrt{2\pi}}\int_{-\infty}^{\infty}g(t)h_1(t)dt\cr
\noalign{\noindent (karena $g$ mewakili transformasi Fourier dari $f$)}
&=\Bover1{\sqrt{2\pi}}\int_{-\infty}^{\infty}g(t)\varhat{h}(y-t)dt
=\Bover1{\sqrt{2\pi}}(g*\varhat{h})(y).\cr}$$

\noindent Karena $y$ sebarang,
$\Bover1{\sqrt{2\pi}}g*\varhat{h}$ adalah transformasi Fourier dari $f\times
h$.


\medskip

{\bf (b)} Tulis $f_1$ untuk transformasi Fourier dari $g\times\varhat{h}$,
$\Reverse{f}(x)=f(-x)$ kapan pun ruas ini terdefinisi, dan $\Reverse{h}(x)=h(-x)$
untuk setiap $x$, sehingga $\Reverse{f}$ mewakili transformasi Fourier dari
$g$, menurut 284If, dan $\Reverse{h}$ adalah transformasi Fourier dari
$\varhat{h}$. Dengan (a), kita memiliki
$f_1=\Bover1{\sqrt{2\pi}}\Reverse{f}*\Reverse{h}$. Ini berarti bahwa
transformasi Fourier invers dari $\sqrt{2\pi}g\times\varhat{h}$ harus
$\sqrt{2\pi}\Reverse{f}_1=(\Reverse{f}*\Reverse{h})\ssplrarrow$; dan sebagai

$$\eqalign{(\Reverse{f}*\Reverse{h})\ssplrarrow(y)
&=(\Reverse{f}*\Reverse{h})(-y)
=\int_{-\infty}^{\infty}\Reverse{f}(t)\Reverse{h}(-y-t)dt\cr
&=\int_{-\infty}^{\infty}f(-t)h(y+t)dt
=\int_{-\infty}^{\infty}f(t)h(y-t)dt
=(f*h)(y),\cr}$$

\noindent transformasi Fourier invers dari $\sqrt{2\pi}g\times\varhat{h}$ adalah
$f*h$ (yang karena itu kontinu), dan $\sqrt{2\pi}g\times\varhat{h}$ harus
mewakili transformasi Fourier dari $f*h$. }%end of proof of 284K

\cmmnt{\medskip

\noindent{\bf Catatan} Bandingkan 283M. Ciri khas teori transformasi Fourier
adalah adanya rumus-rumus yang berlaku dalam berbagai konteks, yang masing-masing
memerlukan interpretasi dan bukti berbeda. }%end of comment

\leader{284L}{}\cmmnt{ Kita sekarang siap untuk hasil yang sepadan dengan
282H. Saya menggunakan metode berbeda, atau setidaknya susunan gagasan yang
berbeda, melalui fakta berikut, yang juga penting untuk keperluan lain.

\medskip

\noindent}{\bf Proposisi} Misalkan $f$ adalah fungsi tempered. Dengan menulis
$\psi_{\sigma}(x) =\bover1{\sigma\sqrt{2\pi}}e^{-x^2/2\sigma^2}$ untuk
$x\in\Bbb R$ dan $\sigma>0$, maka

\Centerline{$\lim_{\sigma\downarrow 0}(f*\psi_{\sigma})(x)=c$}

\noindent setiap kali $x\in\Bbb R$ dan $c\in\Bbb C$ sedemikian sehingga

\Centerline{$\lim_{\delta\downarrow 0}
\Bover1{\delta}\int_{0}^{\delta}|f(x+t)+f(x-t)-2c|dt=0$.}

\proof{{\bf (a) } Menurut 284Bf, setiap $\psi_{\sigma}$ adalah fungsi uji yang
menurun cepat, sehingga menurut 284J $f*\psi_{\sigma}$ terdefinisi di setiap
titik. Kita memerlukan fakta bahwa
$\int_{-\infty}^{\infty}\psi_{\sigma}=1$; dengan substitusi $u=x/\sigma$,

\Centerline{$\int_{-\infty}^{\infty}\psi_{\sigma}
=\Bover1{\sqrt{2\pi}}\int_{-\infty}^{\infty}e^{-u^2/2}du=1$,}

\noindent menurut 263G. Argumen selanjutnya mengikuti pola 282H. Tetapkan

\Centerline{$\phi(t)=|f(x+t)+f(x-t)-2c|$}

\noindent kapan pun ruas ini terdefinisi, yaitu hampir di mana-mana, dan
$\Phi(t)=\int_0^t\phi$, yang terdefinisi untuk semua $t\ge 0$ karena $f$
terintegralkan pada setiap interval terbatas (284Eb). Kita mempunyai

$$\eqalignno{|(f*\psi_{\sigma})(x)-c|
&=|\int_{-\infty}^{\infty}f(x-t)\psi_{\sigma}(t)dt
-c\int_{-\infty}^{\infty}\psi_{\sigma}(t)dt|\cr
&=|\int_{-\infty}^0(f(x-t)-c)\psi_{\sigma}(t)dt
+\int_{0}^{\infty}(f(x-t)-c)\psi_{\sigma}(t)dt|\cr
&=|\int_{0}^{\infty}(f(x+t)-c)\psi_{\sigma}(t)dt
+\int_{0}^{\infty}(f(x-t)-c)\psi_{\sigma}(t)dt|\cr
\noalign{\noindent (karena $\psi_{\sigma}$ adalah fungsi genap)}
&=|\int_{0}^{\infty}(f(x+t)+f(x-t)-2c)\psi_{\sigma}(t)dt|\cr
&\le\int_{0}^{\infty}|f(x+t)+f(x-t)-2c|\psi_{\sigma}(t)dt
=\int_0^{\infty}\phi\times\psi_{\sigma}.\cr}$$

\medskip

{\bf (b)} Kita harus menjelaskan mengapa integral terakhir berhingga. Karena
$f$ adalah fungsi tempered, demikian pula fungsi $t\mapsto f(x+t)$ dan
$t\mapsto f(x-t)$ (284Ec). Fungsi konstan tentu saja tempered, sehingga
$t\mapsto\phi(t)=|f(x+t)+f(x-t)-2c|$ juga tempered. Karena $\psi_{\sigma}$
adalah fungsi uji yang menurun cepat, 284F menunjukkan bahwa hasil kalinya
terintegralkan.

\medskip

{\bf (c)} Ambil $\epsilon>0$. Menurut hipotesis,
$\lim_{t\downarrow 0}\Phi(t)/t=0$; pilih $\delta>0$ sedemikian sehingga
$\Phi(t)\le\epsilon t$ untuk setiap $t\in[0,\delta]$. Ambil sebarang
$\sigma\in\ocint{0,\delta}$. Kita pecah integral
$\int_0^{\infty}\phi\times\psi_{\sigma}$ menjadi tiga bagian.

\medskip

\quad{\bf (i)} Untuk integral dari $0$ sampai $\sigma$, kita mempunyai

$$\int_0^{\sigma}\phi\times\psi_{\sigma}
\le\int_0^{\sigma}\Bover{1}{\sigma\sqrt{2\pi}}\phi
=\Bover1{\sigma\sqrt{2\pi}}\Phi(\sigma)
\le\Bover{\epsilon\sigma}{\sigma\sqrt{2\pi}}\le\epsilon,$$

\noindent karena $\psi_{\sigma}(t)\le\bover1{\sigma\sqrt{2\pi}}$ untuk setiap
$t$.

\medskip

\quad{\bf (ii)} Untuk integral dari $\sigma$ sampai $\delta$, kita mempunyai

$$\eqalignno{\int_{\sigma}^{\delta}\phi\times\psi_{\sigma}
&\le\Bover1{\sigma\sqrt{2\pi}}\int_{\sigma}^{\delta}
  \phi(t)\Bover{2\sigma^2}{t^2}dt\cr
\noalign{\noindent (karena $e^{-t^2/2\sigma^2}
=1/e^{t^2/2\sigma^2}\le
1/(t^2/2\sigma^2)=2\sigma^2/t^2$ untuk setiap $t\ne 0$)}
&=\sigma\sqrt{\Bover2{\pi}}\int_{\sigma}^{\delta}
  \Bover{\phi(t)}{t^2}dt
=\sigma\sqrt{\Bover2{\pi}}\bigl(\Bover{\Phi(\delta)}{\delta^2}
   -\Bover{\Phi(\sigma)}{\sigma^2}
   +\int_{\sigma}^{\delta}\Bover{2\Phi(t)}{t^3}dt\bigr)\cr
\noalign{\noindent (integrasi parsial -- lihat 225F)}
&\le\sigma\bigl(\Bover{\epsilon}{\delta}
   +\int_{\sigma}^{\delta}\Bover{2\epsilon}{t^2}dt\bigr)\cr
\noalign{\noindent (karena $\Phi(t)\le\epsilon t$ untuk $0\le
t\le\delta$ dan $\sqrt{2/\pi}\le 1$)}
&\le\sigma\bigl(\Bover{\epsilon}{\delta}
   +\Bover{2\epsilon}{\sigma}\bigr)
\le 3\epsilon.\cr}$$

\medskip

\quad{\bf (iii)} Untuk integral dari $\delta$ sampai $\infty$, kita mempunyai

$$\eqalign{\int_{\delta}^{\infty}\phi\times\psi_{\sigma}
&=\bover1{\sqrt{2\pi}}\int_{\delta}^{\infty}
  \phi(t)\bover{e^{-t^2/2\sigma^2}}{\sigma}dt.\cr}$$

\noindent Sekarang, untuk setiap $t\ge\delta$,

\Centerline{$\sigma\mapsto\Bover1{\sigma}
e^{-t^2/2\sigma^2}:\ocint{0,\delta}\to\Bbb R$}

\noindent meningkat secara monoton, karena turunan

\Centerline{$\Bover{d}{d\sigma}\Bover1{\sigma}e^{-t^2/2\sigma^2}
=\Bover1{\sigma^2}\bigl(\Bover{t^2}{\sigma^2}-1\bigr) e^{-t^2/2\sigma^2}$}

\noindent adalah positif, dan

\Centerline{$\lim_{\sigma\downarrow 0}\Bover1{\sigma} e^{-t^2/2\sigma^2}=
\lim_{a\to\infty}ae^{-a^2t^2/2}=0$.}

\noindent Jadi kita dapat menerapkan teorema konvergensi terdominasi Lebesgue
untuk memperoleh

$$\lim_{n\to\infty}\int_{\delta}^{\infty}
   \phi(t)\bover{e^{-t^2/2\sigma_n^2}}{\sigma_n}dt=0$$

\noindent untuk setiap barisan $\sequencen{\sigma_n}$ dalam
$\ocint{0,\delta}$ yang konvergen ke $0$, sehingga

$$\lim_{\sigma\downarrow 0}\int_{\delta}^{\infty}
   \phi(t)\bover{e^{-t^2/2\sigma^2}}{\sigma}dt=0.$$

\noindent Oleh karena itu, ada $\sigma_0\in\ocint{0,\delta}$ sedemikian
sehingga

\Centerline{$\int_{\delta}^{\infty}\phi\times\psi_{\sigma} \le\epsilon$}

\noindent untuk setiap $\sigma\le \sigma_0$.

\medskip

\quad{\bf (iv)} Dengan menggabungkan ketiga taksiran tersebut, kita memperoleh

\Centerline{$|(f*\psi_{\sigma})(x)-c|
\le\int_0^{\infty}\phi\times\psi_{\sigma}
\le\epsilon+3\epsilon+\epsilon=5\epsilon$}

\noindent setiap kali $0<\sigma\le\sigma_0$. Karena $\epsilon$ sebarang,
$\lim_{\sigma\downarrow 0}(f*\psi_{\sigma})(x)=c$, sebagaimana diklaim. }%end of proof of 284L

\vleader{60pt}{284M}{Teorema} Misalkan $f$ dan $g$ adalah fungsi-fungsi
tempered sedemikian sehingga $g$ mewakili transformasi Fourier dari $f$. Maka

(a)(i) $g(y) =\lim_{\epsilon\downarrow
0}\Bover1{\sqrt{2\pi}}\int_{-\infty}^{\infty} e^{-iyx}e^{-\epsilon x^2}f(x)dx$
untuk hampir setiap $y\in\Bbb R$.

\quad(ii) Jika $y\in\Bbb R$ sedemikian sehingga $a=\lim_{t\in\dom
g,t\uparrow y}g(t)$ dan $b=\lim_{t\in\dom g,t\downarrow y}g(t)$ keduanya
terdefinisi dalam $\Bbb C$, maka

\Centerline{$\lim_{\epsilon\downarrow 0}\Bover1{\sqrt{2\pi}}
\int_{-\infty}^{\infty} e^{-iyx}e^{-\epsilon x^2}f(x)dx=\Bover12(a+b)$.}

(b)(i) $f(x) =\lim_{\epsilon\downarrow
0}\Bover1{\sqrt{2\pi}}\int_{-\infty}^{\infty} e^{ixy}e^{-\epsilon y^2}g(y)dy$
untuk hampir setiap $x\in\Bbb R$.

\quad(ii) Jika $x\in\Bbb R$ sedemikian sehingga $a=\lim_{t\in\dom
f,t\uparrow x}f(t)$ dan $b=\lim_{t\in\dom f,t\downarrow x}f(t)$ keduanya
terdefinisi dalam $\Bbb C$, maka

\Centerline{$\lim_{\epsilon\downarrow 0}\Bover1{\sqrt{2\pi}}
\int_{-\infty}^{\infty} e^{ixy}e^{-\epsilon y^2}g(y)dy=\Bover12(a+b)$.}

\proof{{\bf (a)(i)} Dengan 223D,

\Centerline{$\lim_{\delta\downarrow 0}\Bover1{2\delta}
\int_{-\delta}^{\delta}|g(y+t)-g(y)|dt=0$}

\noindent untuk hampir setiap $y\in\Bbb R$, karena $g$ terintegralkan pada
setiap interval terbatas. Tetapkan suatu $y$ dengan sifat tersebut. Definisikan
$\phi(t)=|g(y+t)+g(y-t)-2g(y)|$ kapan pun ruas ini terdefinisi. Seperti dalam
bukti 282Ia,

\Centerline{$\int_0^{\delta}\phi \le\int_{-\delta}^{\delta}|g(y+t)-g(y)|dt$,}

\noindent jadi $\lim_{\delta\downarrow
0}\bover1{\delta}\int_0^{\delta}\phi=0$. Akibatnya, menurut 284L,

\Centerline{$g(y)=\lim_{\sigma\to\infty}(g*\psi_{1/\sigma})(y)$.}

\noindent Kita tahu dari 283N bahwa transformasi Fourier dari $\psi_{\sigma}$
adalah $\bover1{\sigma}\psi_{1/\sigma}$ untuk setiap $\sigma>0$. Oleh karena
itu, menurut 284K, $g*\psi_{1/\sigma}$ adalah transformasi Fourier dari
$\sigma\sqrt{2\pi} f\times\psi_{\sigma}$, yaitu,

\Centerline{$(g*\psi_{1/\sigma})(y)
=\int_{-\infty}^{\infty}e^{-iyx}\sigma\psi_{\sigma}(x)f(x)dx$.}

\noindent Jadi,

$$\eqalign{g(y)
&=\lim_{\sigma\to\infty}
   \int_{-\infty}^{\infty}e^{-iyx}\sigma\psi_{\sigma}(x)f(x)dx\cr
&=\lim_{\sigma\to\infty}\Bover1{\sqrt{2\pi}}\int_{-\infty}^{\infty}
   e^{-iyx}e^{-x^2/2\sigma^2}f(x)dx\cr
&=\lim_{\epsilon\downarrow 0}\Bover1{\sqrt{2\pi}}\int_{-\infty}^{\infty}
   e^{-iyx}e^{-\epsilon x^2}f(x)dx.\cr}$$

\noindent dan kesamaan ini berlaku untuk hampir setiap $y$.

\medskip

\quad{\bf (ii) } Sekali lagi, tetapkan $c=\bover12(a+b)$,
$\phi(t)=|g(y+t)+g(y-t)-2c|$ kapan pun ruas ini terdefinisi; kita mempunyai
$\lim_{t\in\dom\phi,t\downarrow 0}\phi(t)\penalty-100=0$, jadi tentu saja
$\lim_{\delta\downarrow 0}\bover{1}{\delta}\int_0^{\delta}\phi=0$, dan

$$c
=\lim_{\sigma\to\infty}(g*\psi_{1/\sigma})(y)
=\lim_{\epsilon\downarrow 0}\Bover1{\sqrt{2\pi}}\int_{-\infty}^{\infty}
   e^{-iyx}e^{-\epsilon x^2}f(x)dx$$

\noindent seperti sebelumnya.


\medskip

{\bf (b)} Pernyataan ini dapat dibuktikan dengan argumen serupa; atau dapat
diturunkan langsung dari (a), dengan mengamati bahwa
$x\mapsto\Reverse{f}(x)=f(-x)$ mewakili transformasi Fourier dari $g$ (lihat
284Ie), lalu menerapkan (a) kepada $g$ dan $\Reverse{f}$. }%end of proof of 284M

\leader{284N}{$L^2$ \dvrocolon{ruang}}\cmmnt{ Kita sekarang siap untuk
hasil-hasil yang sepadan dengan 282J-282K.

\medskip

\noindent}{\bf Lema} Misalkan $\eusm L_{\Bbb C}^2$ adalah ruang fungsi
bernilai kompleks yang terintegralkan-kuadrat pada $\Bbb R$, dan $\eusm S$
ruang fungsi uji yang menurun cepat. Maka untuk setiap
$f\in\eusm L_{\Bbb C}^2$ dan $\epsilon>0$ terdapat $h\in\eusm S$ sedemikian
sehingga $\|f-h\|_2\le\epsilon$.

\proof{ Tetapkan $\phi(x)=e^{-1/x}$ untuk $x>0$ dan nol untuk $x\le 0$;
ingat dari bukti 284G bahwa $\phi$ mulus. Untuk setiap $a<b$, fungsi-fungsi

\Centerline{$x\mapsto\phi_n(x)=\phi(n(x-a))\phi(n(b-x))$}

\noindent membentuk barisan fungsi uji yang konvergen naik ke
$\chi\ooint{a,b}$, sehingga (seperti dalam 284G)

\Centerline{$\inf_{h\in\eusm S}\,\|\chi\ooint{a,b}-h\|^2_2
\le\lim_{n\to\infty}\int_a^b|1-\phi_n|^2 =0$.}

\noindent Karena $\eusm S$ adalah ruang linear (284Ba), untuk setiap fungsi
tangga $g$ dengan dukungan terbatas dan setiap $\epsilon>0$ terdapat
$h\in\eusm S$ sedemikian sehingga $\|g-h\|_2\le\bover12\epsilon$. Namun,
menurut 244H/244Pb, untuk setiap $f\in\eusm L_{\Bbb C}^2$ dan $\epsilon>0$
terdapat fungsi tangga $g$ dengan dukungan terbatas sedemikian sehingga
$\|f-g\|_2\le\bover12\epsilon$; jadi terdapat $h\in\eusm S$ sedemikian sehingga

\Centerline{$\|f-h\|_2\le\|f-g\|_2+\|g-h\|_2\le \epsilon$.}

\noindent Karena $f$ dan $\epsilon$ sebarang, hasilnya terbukti. }%end of proof of 284N

\leader{284O}{Teorema} (a) Misalkan $f$ adalah sebarang fungsi bernilai
kompleks yang terintegralkan-kuadrat pada $\Bbb R$. Maka $f$ adalah fungsi
tempered dan transformasi Fouriernya diwakili oleh fungsi
terintegralkan-kuadrat lain $g$, dengan $\|g\|_2=\|f\|_2$.

(b) Jika $f_1$ dan $f_2$ adalah fungsi bernilai kompleks yang
terintegralkan-kuadrat pada $\Bbb R$, dan transformasi Fourier mereka diwakili
oleh fungsi $g_1$, $g_2$, maka

\Centerline{$\int_{-\infty}^{\infty}f_1\times\bar f_2
=\int_{-\infty}^{\infty}g_1\times\bar g_2$.}

(c) Jika $f_1$ dan $f_2$ adalah fungsi bernilai kompleks yang
terintegralkan-kuadrat pada $\Bbb R$, dan transformasi Fourier mereka diwakili
oleh fungsi $g_1$, $g_2$, maka fungsi terintegralkan $f_1\times f_2$ mempunyai
transformasi Fourier $\bover1{\sqrt{2\pi}}g_1*g_2$.

(d) Jika $f_1$ dan $f_2$ adalah fungsi bernilai kompleks yang
terintegralkan-kuadrat pada $\Bbb R$, dan transformasi Fourier mereka diwakili
oleh fungsi $g_1$, $g_2$, maka $\sqrt{2\pi}g_1\times g_2$ mewakili
transformasi Fourier dari fungsi kontinu $f_1*f_2$.

\proof{{\bf (a)(i)} Pertama, perhatikan kasus ketika $f$ adalah fungsi uji
yang menurun cepat dan $g$ adalah transformasi Fouriernya. Kita tahu bahwa $g$
juga merupakan \rdtf, dan $f$ adalah \Ft\ invers dari $g$ (284C). Konjugat
kompleks $\overline{g}$ dari $g$ diberikan oleh rumus

\Centerline{$\overline{g}(y) =\overline{\Bover1{\sqrt{2\pi}}
\intop\nolimits_{-\infty}^{\infty}e^{-iyx}f(x)dx}
=\Bover1{\sqrt{2\pi}}\int_{-\infty}^{\infty}e^{iyx}\overline{f}(x)dx$,}

\noindent sehingga $\overline{g}$ adalah transformasi Fourier invers dari
$\overline{f}$.

\Centerline{$\int f\times\overline{f} =\int\varcheck{g}\times\overline{f}
=\int g\times\varcheck{\overline{f}} =\int g\times\overline{g}$,}

\noindent dengan menggunakan 283O untuk persamaan tengah.

\medskip

\quad{\bf (ii)} Sekarang misalkan $f\in\eusm L_{\Bbb C}^2$. Tadi telah
dinyatakan bahwa $f$ adalah fungsi tempered; hal ini semata-mata karena

\Centerline{$\int_{-\infty}^{\infty}
\bigl(\Bover{1}{1+|x|}\bigr)^2dx<\infty$,}

\noindent jadi

\Centerline{$\int_{-\infty}^{\infty} \Bover{|f(x)|}{1+|x|}dx<\infty$}

\noindent (244Eb). Dengan 284N, ada barisan $\sequencen{f_n}$ dari
fungsi-fungsi uji yang menurun cepat sedemikian rupa sehingga
$\lim_{n\to\infty}\|f-f_n\|_2=0$. Dengan (i),

\Centerline{$\lim_{m,n\to\infty}\|\varhatf_m-\varhatf_n\|_2
=\lim_{m,n\to\infty}\|f_m-f_n\|_2=0$,}

\noindent dan barisan $\sequencen{{\varhatf}\vthsp_n^{\ssbullet}}$ dari
kelas-kelas ekuivalensi merupakan barisan Cauchy dalam $L_{\Bbb C}^2$. Karena
$L_{\Bbb C}^2$ lengkap (244G/244Pb),
$\sequencen{{\varhatf}\vthsp_n^{\ssbullet}}$ memiliki batas dalam $L_{\Bbb
C}^2$, yang dapat diwakili sebagai $g^{\ssbullet}$ untuk suatu $g\in\eusm
L_{\Bbb C}^2$. Seperti halnya $f$, $g$ harus merupakan fungsi tempered. Tentu saja

\Centerline{$\|g\|_2=\lim_{n\to\infty}\|\varhatf_n\|_2
=\lim_{n\to\infty}\|f_n\|_2=\|f\|_2$.}

\noindent Sekarang, jika $h$ adalah sebarang fungsi uji yang menurun cepat,
$h$ dan $\varhat{h}$ terintegralkan-kuadrat (284Bc, 284C), sehingga

\Centerline{$\int g\times h =\lim_{n\to\infty}\int\varhatf_n\times h
=\lim_{n\to\infty}\int f_n\times\varhat{h} =\int f\times\varhat{h}$.}

\noindent Jadi $g$ mewakili transformasi Fourier dari $f$.

\medskip

{\bf (b)} Menurut 284Ib, setiap fungsi yang mewakili transformasi Fourier dari
$f_1$ atau $f_2$ harus sama hampir di mana-mana dengan fungsi
terintegralkan-kuadrat, sehingga fungsi itu sendiri terintegralkan-kuadrat
dengan norma yang tepat. Seperti dalam 282K (bagian (d) dari buktinya), jika
$g_1$, $g_2$ mewakili transformasi Fourier dari $f_1$, $f_2$, maka
$ag_1+bg_2$ mewakili transformasi Fourier dari $af_1+bf_2$ dan
$\|ag_1+bg_2\|_2=\|af_1+bf_2\|_2$ untuk semua $a$, $b\in\Bbb C$; oleh karena
itu,

\Centerline{$\int f_1\times\overline{f}_2=(f_1|f_2) =(g_1|g_2)=\int
g_1\times\overline{g}_2$.}

\medskip

{\bf (c)} Tentu saja $f_1\times f_2$ terintegralkan karena merupakan hasil
kali dua fungsi terintegralkan-kuadrat (244E/244Pb).

\medskip

\quad{\bf (i)} Ambil $y\in\Bbb R$ dan tetapkan
$f(x)=\overline{f_2(x)}e^{iyx}$ untuk $x\in\Bbb R$. Kemudian $f\in\eusm
L_{\Bbb C}^2$. Kita perlu tahu bahwa transformasi Fourier dari $f$ diwakili
oleh $g$, dengan $g(u)=\overline{g_2(y-u)}$. \Prf\ Misalkan $h$ adalah fungsi
uji yang menurun cepat. Maka

$$\eqalign{\int g\times h
&=\int\overline{g_2(y-u)}h(u)du
=\int \overline{g_2(u)}h(y-u)du\cr
&=\overline{\int g_2\times h_1}
=\overline{\int f_2\times \varhat{h}_1},\cr}$$

\noindent dengan $h_1(u)=\overline{h(y-u)}$. Untuk menghitung $\varhat{h}_1$,
kita mempunyai

$$\eqalign{\varhat{h}_1(v)
&=\Bover1{\sqrt{2\pi}}\int_{-\infty}^{\infty}e^{-ivu}h_1(u)du
=\Bover1{\sqrt{2\pi}}\int_{-\infty}^{\infty}
  e^{-ivu}\overline{h(y-u)}du\cr
&=\overline{\Bover1{\sqrt{2\pi}}\int_{-\infty}^{\infty}
  e^{ivu}h(y-u)du}
=\overline{\Bover1{\sqrt{2\pi}}\int_{-\infty}^{\infty}
  e^{iv(y-u)}h(u)du}
=\overline{e^{ivy}\varhat{h}(v)}.\cr}$$

\noindent Jadi,

\Centerline{$\int g\times h =\overline{\intop f_2\times\varhat{h}_1}
=\int\overline{f_2(v)\varhat{h}_1(v)}dv
=\int\overline{f_2(v)}e^{ivy}\varhat{h}(v)dv =\int f\times\varhat{h}$;}

\noindent karena $h$ sebarang, $g$ mewakili transformasi Fourier
dari $f$. \Qed

\medskip

\quad{\bf (ii)} Sekarang kita punya

$$\eqalignno{(f_1\times f_2)\varsphat(y)
&=\Bover1{\sqrt{2\pi}}
  \int_{-\infty}^{\infty}e^{-iyx}f_1(x)f_2(x)dx\cr
&=\Bover1{\sqrt{2\pi}}
  \int_{-\infty}^{\infty}f_1\times\bar f
=\Bover1{\sqrt{2\pi}}
  \int_{-\infty}^{\infty}g_1\times\bar g\cr
\noalign{\noindent (menggunakan bagian (b))}
&=\Bover1{\sqrt{2\pi}}\int_{-\infty}^{\infty}g_1(u)g_2(y-u)du
=\Bover1{\sqrt{2\pi}}(g_1*g_2)(y).\cr}$$

\noindent Karena $y$ sebarang, $(f_1\times
f_2)\varsphat=\bover1{\sqrt{2\pi}}g_1*g_2$, seperti yang diklaim.

\medskip

{\bf (d)} Menurut (c), transformasi Fourier dari $\sqrt{2\pi}g_1\times g_2$
adalah $\Reverse{f}_1*\Reverse{f}_2$, dengan
$\Reverse{f}_1(x)=f_1(-x)$, sebab $\Reverse{f}_1$ mewakili transformasi Fourier
dari $g_1$. Jadi transformasi Fourier invers dari $\sqrt{2\pi}g_1\times g_2$
adalah $(\Reverse{f}_1*\Reverse{f}_2)\ssplrarrow$. Namun, seperti dalam bukti
284Kb, $(\Reverse{f}_1*\Reverse{f}_2)\ssplrarrow=f_1*f_2$, sehingga $f_1*f_2$
adalah transformasi Fourier invers dari $\sqrt{2\pi}g_1\times g_2$, dan
$\sqrt{2\pi}g_1\times g_2$ mewakili transformasi Fourier dari $f_1*f_2$,
sebagaimana diklaim. Selain itu, karena merupakan transformasi Fourier dari
fungsi terintegralkan, $f_1*f_2$ kontinu (283Cf; lihat juga 255K). } %end of proof of 284O

\leader{284P}{Akibat} Dengan $L_{\Bbb C}^2$ menyatakan ruang Hilbert
kelas-kelas ekuivalensi fungsi bernilai kompleks yang terintegralkan-kuadrat
pada $\Bbb R$, terdapat isometri linear $T:L_{\Bbb C}^2\to L_{\Bbb C}^2$
yang diberikan oleh
$T(f^{\ssbullet})=g^{\ssbullet}$ setiap kali $f$, $g\in\eusm L_{\Bbb C}^2$ dan
$g$ mewakili transformasi Fourier dari $f$.

\cmmnt{ \leader{284Q}{Catatan (a)} Hasil 284P tentu sepadan dengan 282K,
yang menguraikan isometri serupa antara $\ell_{\Bbb C}^2(\Bbb Z)$ dan
$L_{\Bbb C}^2(\ocint{-\pi,\pi})$. %terms reversed for sake of line break
Dalam kasus itu terdapat asimetri nyata yang tidak muncul di sini. Ukuran yang
relevan pada $\Bbb Z$, yaitu ukuran cacah, memberi setiap titik massa tak nol,
sehingga anggota $\ell_{\Bbb C}^2$ merupakan fungsi sejati; tidak mengherankan
bahwa tersedia rumus langsung untuk
$S(f^{\ssbullet})\in\ell_{\Bbb C}^2$ bagi setiap
$f\in\eusm L_{\Bbb C}^2(\ocint{-\pi,\pi})$. Kesulitan menguraikan
$S^{-1}:\ell_{\Bbb C}^2(\Bbb Z)\to L_{\Bbb C}^2(\ocint{-\pi,\pi})$ sangat
mirip dengan kesulitan menguraikan
$T:L_{\Bbb C}^2(\Bbb R)\to L_{\Bbb C}^2(\Bbb R)$ beserta inversnya. Hasil
284Yg dan 286U-286V menunjukkan betapa dekat kemiripan ini.

\header{284Qb}{\bf (b)} Saya menguraikan bagian (c) dan (d) dari 284O secara
rinci, mungkin lebih rinci daripada perlu, karena keduanya memberi kesempatan
untuk menekankan perbedaan antara `$\sqrt{2\pi}g_1\times g_2$ {\it mewakili}
transformasi Fourier dari $f_1*f_2$' dan
`$\bover1{\sqrt{2\pi}}g_1*g_2$ {\it adalah} transformasi Fourier dari
$f_1\times f_2$'. Fungsi $g_1$ dan $g_2$ sendiri tidak ditentukan secara unik
oleh hipotesis bahwa keduanya mewakili transformasi Fourier dari $f_1$ dan
$f_2$, meskipun kelas ekuivalensinya $g_1^{\ssbullet}$,
$g_2^{\ssbullet}\in L_{\Bbb C}^2$ ditentukan. Jadi hasil kali $g_1\times g_2$
juga tidak ditentukan secara unik sebagai fungsi, meskipun kelas ekuivalensinya
$(g_1\times g_2)^{\ssbullet}=g_1^{\ssbullet}\times g_2^{\ssbullet}$ terdefinisi
dengan baik sebagai anggota $L_{\Bbb C}^1$. Sebaliknya, fungsi kontinu
$g_1*g_2$ tidak berubah jika $g_1$ dan $g_2$ diubah pada himpunan terabaikan,
sehingga terdefinisi dengan baik sebagai fungsi. Karena $f_1\times f_2$
terintegralkan dan mempunyai transformasi Fourier sejati, wajar mengharapkan
$(f_1\times f_2)\varsphat$ sama persis dengan
$\bover1{\sqrt{2\pi}}g_1*g_2$.

Perbedaan antara `merupakan' dan `mewakili' transformasi Fourier ini menggemakan
persoalan dalam 233D mengenai ekspektasi bersyarat. Di sana saya menyebut suatu
ekspektasi bersyarat pada $\Tau$ dari fungsi $f$ sebagai fungsi
$\mu\restr\Tau$-terintegralkan $g$ sedemikian sehingga
$\int_Fg\,d\mu=\int_Ffd\mu$ untuk setiap $F\in\Tau$; intinya, setiap fungsi
$\mu\restr\Tau$-terintegralkan yang sama dengan $g$ hampir di mana-mana juga
merupakan ekspektasi bersyarat dari $f$. Demikian pula di sini, jika $g$
mewakili transformasi Fourier dari $f$, setiap fungsi yang sama dengan $g$
hampir di mana-mana juga mewakili transformasi Fourier dari $f$. Dalam 242J
saya menyarankan penyelesaian
komplikasi tersebut dengan memandang ekspektasi bersyarat sebagai pemetaan
antara ruang-ruang $L^1$, bukan antara ruang-ruang $\eusm L^1$. Secara serupa,
transformasi Fourier dalam 284H dapat dipandang sebagai operator linear yang
terdefinisi pada subruang tertentu dari $L^0(\mu)$.

Untuk ekspektasi bersyarat, saya kira ada alasan kuat untuk mengambil operator
pada ruang $L^1$ sebagai perwujudan sejati gagasan tersebut; hal ini akan saya
kembangkan dalam Bab 36 pada jilid berikutnya. Untuk transformasi Fourier,
argumen itu menurut saya tidak sama kuatnya. Dalam 284R di bawah dan \S285,
kita akan melihat kasus-kasus penting ketika kita ingin berbicara tentang
transformasi Fourier yang tidak dapat diwakili oleh anggota $L^0$; jadi sudut
pandang ini masih baru merupakan tahap perantara.

\spheader 284Qc Tentu saja 284Oc-284Od juga memperlihatkan teknik khas dalam
argumen transformasi Fourier: memperluas dengan kontinuitas suatu hubungan
yang berlaku untuk fungsi-fungsi uji.

\spheader 284Qd Hasil 284Oa adalah suatu versi {\bf teorema Plancherel}. Rumus
$\|f\|_2=\|\varhatf\|_2$ adalah {\bf identitas Parseval}. }%end of comment

\cmmnt{ \leader{284R}{Fungsi delta Dirac} Perhatikan fungsi tempered
$\chi\Bbb R$ yang bernilai konstan $1$. Dalam arti apa, jika memang ada, kita
dapat memberikan transformasi Fourier kepada $\chi\Bbb R$?

Jika kita memeriksa $\int\chi\Bbb R\times\varhat{h}$, seperti yang disarankan
dalam 284H, kita mendapatkan

\Centerline{$\int_{-\infty}^{\infty}\chi\Bbb R\times\varhat{h}
=\int_{-\infty}^{\infty}\varhat{h}
=\sqrt{2\pi}\varhat{h}\varcheck{\phantom{h}}(0) =\sqrt{2\pi}h(0)$}

\noindent untuk setiap fungsi uji yang menurun cepat $h$. Tentu saja tidak ada
{ \it fungsi} $g$ sedemikian sehingga $\int g\times h=\sqrt{2\pi}h(0)$ untuk
setiap fungsi uji yang menurun cepat $h$, sebab (dengan argumen 284G) kita akan
mempunyai $\int_{a}^bg=\sqrt{2\pi}$ setiap kali $a<0<b$, sehingga integral tak
tentu $g$ tidak mungkin kontinu pada $0$. Namun terdapat { \it ukuran} pada
$\Bbb R$ dengan sifat yang tepat, yaitu ukuran Dirac $\delta_0$ yang
terkonsentrasi di $0$; ini adalah ukuran probabilitas Radon (257Xa), dan
$\int h\,d\delta_0=h(0)$ untuk setiap fungsi $h$ yang terdefinisi pada $0$.
Dengan demikian, kita mempunyai

\Centerline{$\int_{-\infty}^{\infty}\chi\Bbb R\times\varhat{h}
=\sqrt{2\pi}\int h\,d\delta_0$}

\noindent untuk setiap fungsi uji yang menurun cepat $h$, sehingga wajar
mengatakan bahwa ukuran $\nu=\sqrt{2\pi}\delta_0$ `mewakili transformasi
Fourier dari $\chi\Bbb R$'.

Pada titik ini kita dengan senang hati mencatat bahwa

\Centerline{$\Bover1{\sqrt{2\pi}}\int e^{ixy}\nu(dy)=1$}

\noindent untuk setiap $x\in\Bbb R$, sehingga $\chi\Bbb R$ dapat disebut
transformasi Fourier invers dari $\nu$.

Jika kita menelaah rumus-rumus Teorema 284M, kita memperoleh gagasan yang
selaras dengan pemasangan $\chi\Bbb R$ dan $\nu$ ini.

\Centerline{ $\Bover1{\sqrt{2\pi}}\int_{-\infty}^{\infty} e^{-iyx}e^{-\epsilon
x^2}\chi\Bbb R(x)dx
=\Bover1{\sqrt{2\pi}}\int_{-\infty}^{\infty}e^{-iyx}e^{-\epsilon x^2}dx
=\Bover1{\sqrt{2\epsilon}}e^{-y^2/4\epsilon}$}

\noindent untuk setiap $y\in\Bbb R$, menurut 283N dengan
$\sigma=1/\sqrt{2\epsilon}$. Jadi

\Centerline{$\lim_{\epsilon\downarrow 0}\Bover1{\sqrt{2\pi}}
\int_{-\infty}^{\infty}e^{-iyx}e^{-\epsilon x^2}\chi\Bbb R(x)dx=0$}

\noindent untuk setiap $y\ne 0$, dan transformasi Fourier dari $\chi\Bbb R$
seharusnya nol di setiap titik kecuali $0$. Di sisi lain, fungsi $y\mapsto
\Bover1{\sqrt{2\epsilon}}e^{-y^2/4\epsilon}$ semua memiliki integral
$\sqrt{2\pi}$, yang makin terkonsentrasi di sekitar $0$ ketika $\epsilon$
menurun menuju $0$; ini juga mengarahkan kita langsung kepada $\nu$, ukuran
yang memberi massa $\sqrt{2\pi}$ pada $0$.

Jadi, dengan mengizinkan ukuran selain fungsi, kita dapat memperluas gagasan
transformasi Fourier. Tentu saja kita dapat melangkah jauh lebih jauh. Jika
$h$ adalah fungsi uji yang menurun cepat, maka (karena
$\varcheck{h}\varhat{\phantom{h}}=h$)

\Centerline{$\int_{-\infty}^{\infty}x\varhat{h}(x)dx =-i\sqrt{2\pi}h'(0)$,}

\noindent sehingga fungsi identitas $x\mapsto x$ dapat diberi operator
$h\mapsto -i\sqrt{2\pi}h'(0)$ sebagai transformasi Fouriernya.

Pada titik ini kita memasuki teori sesungguhnya tentang distribusi
(Schwartz), atau `fungsi tergeneralisasi', dan sebaiknya saya berhenti. `Fungsi
delta Dirac' paling alami dipandang sebagai ukuran $\delta_0$ di atas; sebagai
alternatif, sebagai $\Bover1{\sqrt{2\pi}}\varhat{\chi\Bbb R}$. }%end of comment

\exercises{\leader{284W}{Kasus multidimensi} Seperti dalam \S283, saya
memberikan latihan-latihan yang dirancang untuk menunjukkan jalan menuju
generalisasi berdimensi $r$.

\header{284Wa}{\bf (a)} Suatu {\bf fungsi uji yang menurun cepat} pada $\BbbR^r$
adalah fungsi $h:\BbbR^r\to\Bbb C$ sedemikian sehingga (i) $h$ {\bf
halus}, yaitu, semua turunan parsial berulang

\Centerline{$\Bover{\partial^m h}{\partial \xi_{j_1}\ldots\partial\xi_{j_m}}$}

\noindent terdefinisi dan kontinu di setiap titik dalam $\BbbR^r$; (ii)

\Centerline{$\sup_{x\in\BbbR^r}\|x\|^k|h(x)|<\infty$, \quad $\sup_{x\in\Bbb
R^r}\|x\|^k|\Bover{\partial^m h}{\partial
\xi_{j_1}\ldots\partial\xi_{j_m}}(x)|<\infty$}

\noindent untuk setiap $k\in\Bbb N$, $j_1,\ldots,j_m\le r$. Suatu {\bf fungsi
tempered} pada $\BbbR^r$ adalah fungsi terukur bernilai kompleks $f$, yang
terdefinisi hampir di mana-mana dalam $\BbbR^r$ dan, untuk suatu
$k\in\Bbb N$,

\Centerline{$\int_{\Bbb R^r}\Bover1{1+\|x\|^k}|f(x)|dx<\infty$.}

\noindent Tunjukkan bahwa jika $f$ adalah fungsi tempered pada $\BbbR^r$ dan
$h$ adalah fungsi uji yang menurun cepat pada $\BbbR^r$, maka $f\times h$
terintegralkan.

\header{284Wb}{\bf (b)} Tunjukkan bahwa jika $h$ adalah fungsi uji yang
menurun cepat pada $\BbbR^r$ maka juga $\varhat{h}$, dan bahwa dalam hal ini
$\varhat{h}\varcheck{\phantom{h}}=h$.

\header{284Wc}{\bf (c)} Tunjukkan bahwa jika $f$ adalah fungsi tempered pada
$\Bbb R^r$ dan $\int f\times h=0$ untuk setiap fungsi uji yang menurun cepat
$h$ pada $\BbbR^r$, maka $f=0$ a.e.

\header{284Wd}{\bf (d)} Jika $f$ dan $g$ adalah fungsi-fungsi tempered pada
$\Bbb R^r$, kita mengatakan bahwa {\bf $g$ mewakili transformasi Fourier dari
$f$} jika $\int g\times h=\int f\times\varhat{h}$ untuk setiap fungsi uji yang
menurun cepat $h$ pada $\BbbR^r$. Tunjukkan bahwa jika $f$ terintegralkan,
maka $\varhatf$ mewakili transformasi Fourier dari $f$ dalam
arti ini.

\header{284We}{\bf (e)} Misalkan $f$ adalah fungsi tempered pada $\BbbR^r$.
Dengan menulis
$\psi_{\sigma}(x)=\bover1{(\sigma\sqrt{2\pi})^r}e^{-x\dotproduct
x/2\sigma^2}$ untuk $x\in\BbbR^r$, tunjukkan bahwa
$\lim_{\sigma\downarrow 0}(f*\psi_{\sigma})(x)=c$ setiap kali
$x\in\BbbR^r$, $c\in\Bbb C$ memenuhi
$\lim_{\delta\downarrow 0}\bover1{\delta^r}
\int_{B(x,\delta)}|f(t)-c|dt=0$, dengan
$B(x,\delta)=\{t:\|t-x\|\le\delta\}$.

\header{284Wf}{\bf (f)} Misalkan $f$ dan $g$ adalah fungsi-fungsi tempered pada
$\BbbR^r$, $g$ mewakili transformasi Fourier dari $f$, dan $h$ adalah fungsi
uji yang menurun cepat. Tunjukkan bahwa (i) transformasi Fourier dari
$f\times h$ adalah $\bover1{(\sqrt{2\pi})^r}g*\varhat{h}$ (ii)
$(\sqrt{2\pi})^rg\times\varhat{h}$ mewakili transformasi Fourier dari $f*h$.

\header{284Wg}{\bf (g)} Misalkan $f$ dan $g$ adalah fungsi-fungsi tempered pada
$\BbbR^r$ sedemikian sehingga $g$ mewakili transformasi Fourier dari $f$.
Tunjukkan bahwa

\Centerline{$g(y) =\lim_{\epsilon\downarrow 0}\Bover1{(\sqrt{2\pi})^r}
\int_{\Bbb R^r} e^{-iy\dotproduct x}e^{-\epsilon x\dotproduct x}f(x)dx$}

\noindent untuk hampir setiap $y\in\BbbR^r$.

\header{284Wh}{\bf (h)} Tunjukkan bahwa untuk setiap fungsi bernilai kompleks
$f$ yang terintegralkan-kuadrat pada $\BbbR^r$ dan setiap $\epsilon>0$ terdapat
\rdtf\ $h$ sedemikian sehingga $\|f-h\|_2\le\epsilon$.

\header{284Wi}{\bf (i)} Misalkan $\eusm L_{\Bbb C}^2$ adalah ruang fungsi
bernilai kompleks yang terintegralkan-kuadrat pada $\BbbR^r$.
Tunjukkan bahwa

\quad (i) untuk setiap $f\in\eusm L_{\Bbb C}^2$ ada $g\in\eusm L_{\Bbb C}^2$
yang mewakili transformasi Fourier dari $f$, dan dalam hal ini
$\|g\|_2=\|f\|_2$;

\quad (ii) jika $g_1$, $g_2\in\eusm L_{\Bbb C}^2$ mewakili transformasi
Fourier dari $f_1$, $f_2\in\eusm L_{\Bbb C}^2$, maka
$\bover1{(\sqrt{2\pi})^r}g_1*g_2$ adalah transformasi Fourier dari $f_1\times
f_2$, dan $(\sqrt{2\pi})^rg_1\times g_2$ mewakili transformasi Fourier dari
$f_1*f_2$.

\spheader 284Wj Misalkan $T$ adalah matriks real taksingular berukuran
$r\times r$, yang dipandang sebagai operator linear dari $\BbbR^r$ ke dirinya
sendiri. (i) Tunjukkan bahwa $\varhatf=|\det T|(fT)\varsphat T\trs$ untuk setiap
fungsi bernilai kompleks $f$ yang terintegralkan pada $\BbbR^r$. (ii)
Tunjukkan bahwa $hT$ adalah \rdtf\ untuk setiap \rdtf\ $h$. (iii) Tunjukkan
bahwa jika $f$, $g$ adalah fungsi-fungsi tempered dan $g$ mewakili transformasi
Fourier dari $f$, maka $\Bover1{|\det T|}g(T\trs)^{-1}$ mewakili transformasi
Fourier dari $fT$; sehingga jika $T$ ortogonal, maka $gT$ mewakili transformasi
Fourier dari $fT$.
%mt28bits
}%end of exercises 284W

\exercises{ \leader{284X}{Latihan dasar (a)} Tunjukkan bahwa jika $g$ dan $h$
adalah fungsi-fungsi uji yang menurun cepat, maka $g\times h$ juga demikian. %284A

\spheader 284Xb Tunjukkan bahwa ada fungsi kontinu terintegralkan tak nol
$f$, $g:\Bbb R\to\Bbb C$ sedemikian sehingga $f*g=0$ di setiap titik.
\Hint{Ambil keduanya sebagai \Ft\ dari fungsi-fungsi uji yang sesuai.}
%284C

\spheader 284Xc Misalkan $f:\Bbb R\to\Bbb C$ adalah fungsi terdiferensialkan
sedemikian sehingga turunannya $f'$ adalah fungsi tempered dan, untuk suatu
$k\in\Bbb N$,

\Centerline{$\lim_{x\to\infty}x^{-k}f(x) =\lim_{x\to-\infty}x^{-k}f(x)=0$.}

\noindent (i) Tunjukkan bahwa $\int f\times h'=-\int f'\times h$ untuk setiap
fungsi uji yang menurun cepat $h$. (ii) Tunjukkan bahwa jika $g$ adalah fungsi
tempered yang mewakili transformasi Fourier dari $f$, maka $y\mapsto iyg(y)$
mewakili transformasi Fourier dari $f'$. %284H

\spheader 284Xd\dvAnew{2013} Untuk fungsi tempered $f$ dan $\alpha\in\Bbb R$,
tetapkan

\Centerline{$(S_{\alpha}f)(x)=f(x+\alpha)$, \quad$(M_{\alpha}f)(x)=e^{i\alpha
x}f(x)$, \quad$(D_{\alpha}f)(x)=f(\alpha x)$}

\noindent kapan pun ruas-ruas tersebut terdefinisi. (i) Tunjukkan bahwa
$S_{\alpha}f$, $M_{\alpha}f$ dan (jika $\alpha\ne 0$) $D_{\alpha}f$ adalah
fungsi-fungsi tempered. (ii) Tunjukkan bahwa jika $g$ adalah fungsi tempered
yang mewakili transformasi Fourier dari $f$, maka $M_{-\alpha}g$ mewakili
transformasi Fourier dari $S_{\alpha}f$, $S_{-\alpha}g$ mewakili transformasi
Fourier dari $M_{\alpha}f$, $\bar{\Reverse{g}}=\Reverse{\bar g}$ mewakili
transformasi Fourier dari $\bar f$, dan jika $\alpha\ne 0$, maka
$\Bover1{|\alpha|}D_{1/\alpha}g$ mewakili transformasi Fourier dari
$D_{\alpha}f$.
%284I

\spheader 284Xe Tunjukkan bahwa jika $h$ adalah fungsi uji yang menurun cepat
dan $f$ adalah fungsi terukur bernilai kompleks, terdefinisi hampir di mana-mana
dalam $\Bbb R$, sedemikian sehingga
$\int_{-\infty}^{\infty}|x|^k|f(x)|dx<\infty$ untuk setiap $k\in\Bbb N$, maka
konvolusi $f*h$ adalah fungsi uji yang menurun cepat.
\Hint{Tunjukkan bahwa transformasi Fourier dari $f*h$ adalah fungsi uji.}
%284J

\sqheader 284Xf Misalkan $f$ adalah fungsi tempered sedemikian sehingga
$\lim_{a\to\infty}\int_{-a}^af$ ada dalam $\Bbb C$. Tunjukkan bahwa batas ini
juga sama dengan $\lim_{\epsilon\downarrow 0}\int_{-\infty}^{\infty}
e^{-\epsilon x^2}f(x)dx$. \Hint{Tetapkan $g(x)=f(x)+f(-x)$. Gunakan 224J untuk
menunjukkan bahwa jika $0\le a\le b$, maka
$|\int_a^bg(x)e^{-\epsilon x^2}dx|\le\sup_{c\in[a,b]}|\int_a^cg|$, sehingga
$\lim_{a\to\infty}\int_0^ag(x)e^{-\epsilon x^2}dx$ ada secara seragam terhadap
$\epsilon$, sementara $\lim_{\epsilon\downarrow 0}\int_0^ag(x)e^{-\epsilon x^2}dx
=\int_0^ag$ untuk setiap $a\ge 0$.} %284L

\sqheader 284Xg Misalkan $f$ dan $g$ adalah fungsi-fungsi tempered pada
$\Bbb R$ sedemikian sehingga $g$ mewakili transformasi Fourier dari $f$.
Tunjukkan bahwa

\Centerline{$g(y)
=\lim_{a\to\infty}\Bover1{\sqrt{2\pi}}\int_{-a}^ae^{-iyx}f(x)dx$}

\noindent pada hampir semua titik $y$ untuk mana batasnya ada. \Hint{284Xf, 284M.} %284M, 284Xf, 284L

\sqheader 284Xh Misalkan $f$ adalah fungsi bernilai kompleks yang terintegralkan
pada $\Bbb R$, dan $\varhatf$ juga terintegralkan. Tunjukkan bahwa
$\varhatf\varcheck{\phantom{f}}=f$ pada setiap titik tempat $f$ kontinu. %284M

\spheader 284Xi Tunjukkan bahwa untuk setiap $p\in\coint{1,\infty}$,
$f\in\eusm L_{\Bbb C}^p$ dan $\epsilon>0$, terdapat fungsi uji yang menurun
cepat $h$ sedemikian sehingga $\|f-h\|_p\le\epsilon$.
%284N

\sqheader 284Xj Misalkan $f$ dan $g$ adalah fungsi-fungsi bernilai kompleks
yang terintegralkan-kuadrat pada $\Bbb R$, dan $g$ mewakili transformasi
Fourier dari $f$. Tunjukkan bahwa

\Centerline{$\int_c^df=\Bover{i}{\sqrt{2\pi}}
\int_{-\infty}^{\infty}\Bover{e^{icy}-e^{idy}}{y}g(y)dy$}

\noindent setiap kali $c<d$ dalam $\Bbb R$. %284O

\spheader 284Xk Misalkan $f$ adalah fungsi terukur bernilai kompleks,
terdefinisi hampir di mana-mana dalam $\Bbb R$, sedemikian sehingga
$\int|f|^p<\infty$, dengan $1<p\le 2$. Tunjukkan bahwa $f$ adalah fungsi
tempered dan ada fungsi tempered $g$ yang mewakili transformasi Fourier dari
$f$. \Hint{Nyatakan $f$ sebagai $f_1+f_2$, dengan $f_1$ terintegralkan dan
$f_2$ terintegralkan-kuadrat.} ({\bf Catatan} Dengan mendefinisikan
$\|f\|_p$, $\|g\|_q$ seperti dalam 244D, dengan $q=p/(p-1)$, kita mempunyai
$\|g\|_q\le(2\pi)^{(p-2)/2p}\|f\|_p$; lihat {\smc Zygmund 59}, XVI.3.2.)
%284O

\spheader 284Xl Misalkan $f$, $g$ adalah fungsi-fungsi bernilai kompleks yang
terintegralkan-kuadrat pada $\Bbb R$, dan $g$ mewakili transformasi Fourier
dari $f$.

\quad{(i)} Tunjukkan bahwa

\Centerline{$\Bover1{\sqrt{2\pi}}\int_{-a}^ae^{ixy}g(y)dy
=\Bover1{\pi}\int_{-\infty}^{\infty}\Bover{\sin at}{t}f(x-t)dt$}

\noindent setiap kali $x\in\Bbb R$ dan $a>0$. \Hint{Temukan transformasi
Fourier invers dari $y\mapsto e^{-ixy}\chi[-a,a](y)$, lalu gunakan 284Ob.}

\quad{(ii)} Tunjukkan bahwa jika $f(x)=0$ untuk $x\in\ooint{c,d}$, maka

\Centerline{$\Bover1{\sqrt{2\pi}}\lim_{a\to\infty}\int_{-a}^a
e^{ixy}g(y)dy=0$}

\noindent untuk $x\in\ooint{c,d}$.

\quad{(iii)} Tunjukkan bahwa jika $f$ dapat dibedakan pada $x\in\Bbb R$, maka

\Centerline{$\Bover1{\sqrt{2\pi}}\lim_{a\to\infty}\int_{-a}^a
e^{ixy}g(y)dy=f(x)$.}

\quad{(iv)} Tunjukkan bahwa jika $f$ mempunyai variasi terbatas pada suatu
interval yang memuat $x$ dalam interiornya, maka

\Centerline{$\Bover1{\sqrt{2\pi}}\lim_{a\to\infty}\int_{-a}^a
e^{ixy}g(y)dy=\Bover12(\lim_{t\in\dom f,t\uparrow x}f(t)+\lim_{t\in\dom
f,t\downarrow x}f(t))$.}
%284O

\spheader 284Xm Misalkan $f$ adalah fungsi kompleks yang terintegralkan pada
$\Bbb R$. Tunjukkan bahwa jika $\varhatf$ terintegralkan-kuadrat, maka $f$ juga
terintegralkan-kuadrat. %284O

\spheader 284Xn Misalkan $f_1$, $f_2$ adalah fungsi-fungsi bernilai kompleks
yang terintegralkan-kuadrat pada $\Bbb R$, dengan \Ft{} masing-masing diwakili
oleh $g_1$, $g_2$. Tunjukkan bahwa
$\int_{-\infty}^{\infty}f_1(t)f_2(-t)dt
=\int_{-\infty}^{\infty}g_1(t)g_2(t)dt$. %284O

\spheader 284Xo Misalkan $x\in\Bbb R$. Tulis $\delta_x$ untuk ukuran Dirac
pada $\Bbb R$ yang terkonsentrasi di $x$. Jelaskan suatu arti yang memungkinkan
$\sqrt{2\pi}\delta_x$ dipandang sebagai transformasi Fourier dari fungsi
$t\mapsto e^{ixt}$.
%284R

\spheader 284Xp Untuk sebarang fungsi tempered $f$ dan $x\in\Bbb R$, misalkan
$\delta_x$ adalah ukuran Dirac pada $\Bbb R$ yang terkonsentrasi di $x$, dan
tetapkan

\Centerline{$(\delta_x*f)(u)=\int f(u-t)\delta_x(dt)=f(u-x)$}

\noindent untuk setiap $u$ sedemikian sehingga $u-x\in\dom f$ (bdk.\ 257Xe).
Jika $g$ mewakili transformasi Fourier dari $f$, carilah representasi yang
sesuai bagi transformasi Fourier dari $\delta_x*f$, dan hubungkan representasi
itu dengan hasil kali $g$ dan transformasi Fourier dari $\delta_x$. %284R, 284Xo

\spheader 284Xq(i) Tunjukkan bahwa

\Centerline{$\lim_{\delta\downarrow 0,a\to\infty}
\bigl(\int_{-a}^{-\delta}\Bover1x e^{-iyx}dx +\int_{\delta}^a\Bover1x
e^{-iyx}dx\bigr) =-\pi i\sgn y$}

\noindent untuk setiap $y\in\Bbb R$, dengan $\sgn y=y/|y|$ jika $y\ne 0$ dan
$\sgn 0=0$. \Hint{283Da.}

\quad{(ii)} Tunjukkan bahwa

\Centerline{$\lim_{c\to\infty}\Bover1c\int_0^c \int_{-a}^ae^{ixy}\sgn
y\,dy\,da =\Bover{2i}x$}

\noindent untuk setiap $x\ne 0$.

\quad{(iii)} Tunjukkan bahwa untuk setiap fungsi uji yang menurun cepat
$h$,

$$\eqalign{\int_0^{\infty}\Bover1x(\varhat{h}(x)-\varhat{h}(-x))dx
&=\lim_{\delta\downarrow 0,a\to\infty}\bigl(\int_{-a}^{-\delta}
\Bover1x\varhat{h}(x)dx
+\int_{\delta}^{a}\Bover1x\varhat{h}(x)dx\bigr)\cr
&=-\Bover{i\pi}{\sqrt{2\pi}}\int_{-\infty}^{\infty}h(y)\sgn
y\,dy.\cr}$$

\quad{(iv)} Tunjukkan bahwa untuk setiap fungsi uji yang menurun cepat $h$,

\Centerline{$\Bover{i\pi}{\sqrt{2\pi}}\int_{-\infty}^{\infty}
  \varhat{h}(x)\sgn x\,dx
=\int_0^{\infty}\Bover1y(h(y)-h(-y))dy$.}
%284notes

\spheader 284Xr Misalkan $\sequencen{h_n}$ adalah barisan fungsi uji yang
menurun cepat sedemikian sehingga
$\phi(f)=\lim_{n\to\infty}\int_{-\infty}^{\infty}h_n\times f$ terdefinisi
untuk setiap fungsi uji yang menurun cepat $f$. Tunjukkan bahwa
$\lim_{n\to\infty}\int_{-\infty}^{\infty}h'_n\times f$,
$\lim_{n\to\infty}\int_{-\infty}^{\infty}\varhat{h}_n\times f$ dan
$\lim_{n\to\infty}\int_{-\infty}^{\infty}(h_n*g)\times f$ terdefinisi untuk
semua fungsi uji yang menurun cepat $f$ dan $g$, dan bernilai nol jika $\phi$
identik nol. \Hint{Hasil 255G akan membantu untuk batas terakhir.} %284+

\leader{284Y}{Latihan lanjutan (a)} %\spheader 284Ya
Misalkan $f$ adalah fungsi bernilai kompleks yang terintegralkan pada
$\ocint{-\pi,\pi}$, dan $\tilde f$ perluasan periodiknya, seperti dalam 282Ae.
Tunjukkan bahwa $\tilde f$ adalah fungsi tempered. Tunjukkan bahwa untuk setiap
fungsi uji yang menurun cepat $h$,
$\int\tilde f\times\varhat{h}
=\sqrt{2\pi}\sum_{k=-\infty}^{\infty}c_kh(k)$, dengan
$\sequence{k}{c_k}$ adalah barisan koefisien Fourier dari $f$. ({\it
Petunjuk\/}: mulai dengan kasus $f(x)=e^{inx}$. Selanjutnya, tunjukkan bahwa

\Centerline{$M=\sum_{k=-\infty}^{\infty}|h(k)| +\sum_{k=-\infty}^{\infty}
\sup_{x\in[(2k-1)\pi,(2k+1)\pi]}|\varhat{h}(x)|<\infty$,}

\noindent dan bahwa

\Centerline{$|\int\tilde f\times\varhat{h}
-\sqrt{2\pi}\sum_{k=-\infty}^{\infty}c_kh(k)| \le M\|f\|_1$.}

\noindent Terakhir, terapkan 282Ib.) %284C

\header{284Yb}{\bf (b)} Misalkan $f$ adalah fungsi bernilai kompleks yang
terdefinisi hampir di mana-mana dalam $\Bbb R$, sedemikian sehingga $f\times h$
terintegralkan untuk setiap fungsi uji yang menurun cepat $h$. Tunjukkan bahwa
$f$ tempered. %284F

\spheader 284Yc Misalkan $f$ dan $g$ adalah fungsi-fungsi tempered pada
$\Bbb R$, dan $g$ mewakili transformasi Fourier dari $f$. Tunjukkan bahwa

$$\int_c^df
=\Bover{i}{\sqrt{2\pi}}\lim_{\sigma\to\infty}\int_{-\infty}^{\infty}
\Bover{e^{icy}-e^{idy}}{y}e^{-y^2/2\sigma^2}g(y)dy$$

\noindent setiap kali $c\le d$ dalam $\Bbb R$. \Hint{Tetapkan
$\theta=\chi[c,d]$. Tunjukkan bahwa kedua ruas sama dengan
$\lim_{\sigma\to\infty}\int f\times(\theta*\psi_{1/\sigma})$, dengan
$\psi_{\sigma}$ didefinisikan seperti dalam 283N dan 284L.} %284H

\spheader 284Yd Tunjukkan bahwa jika $g:\Bbb R\to\Bbb R$ adalah fungsi ganjil
yang mempunyai variasi terbatas dan
$\int_1^{\infty}\bover1x g(x)dx=\infty$, maka $g$ tidak mewakili transformasi
Fourier dari fungsi tempered mana pun. \Hint{283Ye, 284Yc.}
%284Yc, 284H

\ifdim\pagewidth>467pt\fontdimen3\tenrm=1.84pt
 \fontdimen4\tenrm=1.22pt\fi
\spheader 284Ye Misalkan $\eusm S$ adalah ruang fungsi uji yang menurun cepat.
Untuk $k$, $m\in\Bbb N$, tetapkan
$\tau_{km}(h)=\sup_{x\in\Bbb R}|x|^k|h^{(m)}(x)|$ untuk setiap $h\in\eusm S$,
dengan $h^{(m)}$ seperti biasa menyatakan turunan ke-$m$ dari $h$. (i)
Tunjukkan bahwa setiap $\tau_{km}$ adalah seminorma dan bahwa $\eusm S$
lengkap serta separabel untuk topologi ruang linear termetrik $\frak T$ yang
didefinisikan oleh seminorma-seminorma tersebut. (ii) Tunjukkan bahwa
$h\mapsto\varhat{h}:\eusm S\to\eusm S$ kontinu terhadap $\frak T$. (iii)
Tunjukkan bahwa jika $f$ adalah sebarang fungsi tempered, maka
$h\mapsto\int f\times h$ kontinu terhadap $\frak T$. (iv) Tunjukkan bahwa jika
$f$ adalah fungsi terintegralkan sedemikian sehingga
$\int|x^kf(x)|dx<\infty$ untuk setiap $k\in\Bbb N$, maka
$h\mapsto f*h:\eusm S\to\eusm S$ kontinu terhadap $\frak T$. %284J
\fontdimen3\tenrm=1.67pt\fontdimen4\tenrm=1.22pt

\spheader 284Yf Tunjukkan bahwa jika $f$ adalah fungsi tempered pada $\Bbb R$
dan

\Centerline{$\gamma =\lim_{c\to\infty}\Bover1c\int_{0}^c\int_{-a}^af(x)dxda$}

\noindent terdefinisi dalam $\Bbb C$, maka $\gamma$ juga sama dengan

\Centerline{$\lim_{\epsilon\downarrow 0}\int_{-\infty}^{\infty}
f(x)e^{-\epsilon|x|}dx$.}
%284Xf, 284L

\spheader 284Yg Misalkan $f$, $g$ adalah fungsi-fungsi bernilai kompleks yang
terintegralkan-kuadrat pada $\Bbb R$, dan $g$ mewakili transformasi Fourier
dari $f$. Misalkan $m\in\Bbb Z$ dan $(2m-1)\pi<x<(2m+1)\pi$. Tetapkan
$\tilde f(t)=f(t+2m\pi)$ untuk semua $t\in\ocint{-\pi,\pi}$ sedemikian sehingga
$t+2m\pi\in\dom f$. Misalkan $\langle c_k\rangle_{k\in\Bbb Z}$ adalah barisan
koefisien Fourier dari $\tilde f$. Tunjukkan bahwa

\Centerline{$\Bover1{\sqrt{2\pi}}\lim_{a\to\infty}\int_{-a}^a e^{ixy}g(y)dy
=\lim_{n\to\infty}\sum_{k=-n}^nc_ke^{ikx}$}

\noindent dalam arti bahwa jika salah satu batas ada dalam $\Bbb C$, maka batas
yang lain juga ada dan keduanya sama. \Hint{284Xl(i), 282Da.} %284O, 284Xl

\spheader 284Yh Tunjukkan bahwa jika $f$ terintegralkan pada $\Bbb R$ dan ada
$M\ge 0$ sedemikian sehingga $f(x)=\varhatf(x)=0$ untuk $|x|\ge M$, maka
$f=0$ a.e. \Hint{Reduksikan ke kasus $M=\pi$. Dengan menelaah deret Fourier
dari $f\restr\ocint{-\pi,\pi}$, tunjukkan bahwa $f$ dapat dinyatakan sebagai
$f(x)=\sum_{k=-m}^mc_ke^{ikx}$ untuk hampir setiap
$x\in\ocint{-\pi,\pi}$. Kemudian hitung $\varhatf(2n+\bover12)$ untuk $n$
besar.}

\spheader 284Yi Misalkan $\nu$ adalah ukuran Radon pada $\Bbb R$ yang
`tempered', dalam arti bahwa
$\int_{-\infty}^{\infty}\Bover1{1+|x|^k}\nu(dx)$ berhingga untuk suatu
$k\in\Bbb N$. (i) Tunjukkan bahwa setiap fungsi uji yang menurun cepat
$\nu$-terintegralkan. (ii) Tunjukkan bahwa jika $\nu$ mempunyai dukungan
terbatas (definisi: 256Xf), dan $h$ adalah fungsi uji yang menurun cepat, maka
$\nu*h$ adalah fungsi uji yang menurun cepat, dengan
$(\nu*h)(x)=\int_{-\infty}^{\infty}h(x-y)\nu(dy)$ untuk $x\in\Bbb R$. (iii)
Tunjukkan bahwa terdapat barisan $\sequencen{h_n}$ fungsi uji yang menurun
cepat sedemikian sehingga
$\lim_{n\to\infty}\int_{-\infty}^{\infty}h_n\times f
=\int_{-\infty}^{\infty}fd\nu$ untuk setiap fungsi uji yang menurun cepat
$f$. %284Xr.

\spheader 284Yj Misalkan $\phi:\eusm S\to\Bbb C$ adalah fungsional yang
didefinisikan oleh rumus 284Xr. Tunjukkan bahwa $\phi$ kontinu untuk topologi
dalam 284Ye. ({\it Catatan\/}: akan membantu jika pembaca mengetahui sedikit
lebih banyak tentang ruang topologis linear termetrik daripada yang dibahas
dalam \S2A5.) %284Xr

}%end of exercises

\cmmnt{ \Notesheader{284} Sekali lagi saya harus memperingatkan bahwa materi di
atas hanya memberikan pandangan yang sangat terbatas tentang pokok ini. Saya
telah mencoba menunjukkan bagaimana teori transformasi Fourier untuk fungsi
`baik' -- yang di sini berarti fungsi uji yang menurun cepat -- dapat diperluas,
melalui suatu dualitas, ke kelas fungsi yang jauh lebih luas, yaitu
`fungsi-fungsi tempered'. Dengan $\eusm S$ menyatakan ruang linear fungsi uji
yang menurun cepat, kita dapat mencoba menyelidiki transformasi Fourier setiap
fungsional linear $\phi:\eusm S\to\Bbb C$, dengan menulis
$\varhat\phi(h)=\phi(\varhat{h})$ untuk setiap $h\in\eusm S$. (Sebenarnya,
lebih lazim pada tahap ini untuk membatasi perhatian pada fungsional $\phi$
yang kontinu terhadap topologi standar pada $\eusm S$ yang diuraikan dalam
284Ye; fungsional tersebut disebut {\bf distribusi tempered (Schwartz)}.)
Menurut 284F-284G, sebagian fungsional ini dapat diidentifikasi dengan kelas
ekuivalensi fungsi tempered; kemudian kita dapat menyelidiki fungsi-fungsi
tempered yang transformasi Fouriernya kembali dapat diwakili oleh fungsi
tempered.

Struktur teori transformasi Fourier mungkin paling jelas diuraikan melalui
rumus-rumusnya. Tujuan kita ialah membentuk pasangan
$(f,g)=(f,\varhatf)=(\varcheck{g},g)$ sedemikian sehingga berlaku

\quad { \it Inversi\/}: $\varhat{h}\varcheck{\phantom{h}}
=\varcheck{h}\varhat{\phantom{h}}=h$; %284If

\quad { \it Pembalikan\/}: $\varcheck{h}(y)=\varhat{h}(-y)$; %284If

\quad { \it Linearitas\/}: $(h_1+h_2)\varsphat=\varhat{h}_1+\varhat{h}_2$, \quad$(ch)\varsphat=c\varhat{h}$; %284Ic

\quad {\it Diferensiasi\/}: $(h')\varsphat(y)=iy\varhat{h}(y)$;

\quad {\it Pergeseran\/}: jika $h_1(x)=h(x+c)$ maka $\varhat{h}_1(y)=e^{iyc}\varhat{h}(y)$; %284Xd

\quad { \it Modulasi\/}: jika $h_1(x)=e^{icx}h(x)$ maka $\varhat{h}_1(y)=\varhat{h}(y-c)$; %284Xd

\quad {\it Simetri\/}: jika $h_1(x)=h(-x)$ maka $\varhat{h}_1(y)
=\varhat{h}(-y)$; %284Xd

\quad { \it Konjugasi kompleks\/}: $(\overline{h})\varsphat(y)=\overline{\varhat{h}(-y)}$; %284Xd

\quad {\it Dilasi\/}: jika $h_1(x)=h(cx)$, di mana $c>0$, maka $\varhat{h}_1(y)=\Bover1c\varhat{h}(\Bover{y}{c})$; %284Xd

\quad {\it Konvolusi\/}:
$(h_1*h_2)\varsphat=\sqrt{2\pi}\varhat{h}_1\times\varhat{h}_2$, %284Kb
\quad $(h_1\times h_2)\varsphat
=\Bover1{\sqrt{2\pi}}\varhat{h}_1*\varhat{h}_2$; %284Ka

\quad {\it Dualitas\/}: $\int_{-\infty}^{\infty}h_1\times\varhat{h}_2
=\int_{-\infty}^{\infty}\varhat{h}_1\times h_2$;

\quad { \it Parseval\/}: $\int_{-\infty}^{\infty}h_1\times\overline{h_2}
=\int_{-\infty}^{\infty}\varhat{h}_1\times\overline{\varhat{h}_2}$; %284Ob

\noindent Dan, tentu saja,

\Centerline{$\varhat{h}(y)
=\Bover1{\sqrt{2\pi}}\int_{-\infty}^{\infty}e^{-iyx}h(x)dx$,} %284Xg

\Centerline{$\int_c^d\varhat{h}(y)dy
=\Bover{i}{\sqrt{2\pi}}\int_{-\infty}^{\infty}
\Bover{e^{-idy}-e^{-icy}}{y}h(y)dy$.}  %284Xj

\noindent (Saya menggunakan huruf $h$ dalam daftar di atas untuk menandai fakta
bahwa semua rumus tersebut memang berlaku bagi fungsi uji yang menurun cepat.)
Selain itu, sering kali penting bahwa operasi $h\mapsto\varhat{h}$ kontinu
dalam suatu arti.

Tantangan teori transformasi Fourier `murni' ialah menemukan ragam objek $h$
seluas mungkin yang membuat rumus-rumus di atas berlaku, dengan interpretasi
yang tepat bagi $\varsphat$, $*$ dan $\int_{-\infty}^{\infty}$. Perlu dicatat
bahwa sejak awal bidang ini telah diperkaya oleh penerapannya dalam cabang lain
matematika, ilmu fisika dan ilmu sosial. Berulang kali penerapan tersebut
menyarankan pasangan baru $(f,\varhatf)$ dan menuntut kemampuan lebih besar
untuk menafsirkan aturan yang hendak kita ikuti. Bahkan teori distribusi
tampaknya tidak memberi uraian kanonik lengkap tentang semua yang dapat
dilakukan. Pertama, terdapat kesulitan besar dalam menafsirkan `hasil kali' dua
distribusi sebarang, sehingga beberapa rumus di atas bermasalah; kedua, tidak
jelas bahwa cukup mempertimbangkan satu jenis distribusi saja. Dalam bagian ini
saya hanya menelaah satu ruang `fungsi uji', yaitu ruang $\eusm S$ dari fungsi
uji yang menurun cepat; tetapi sedikitnya dua ruang lain juga penting: ruang
$\eusm D$ dari fungsi mulus dengan dukungan terbatas, dan ruang $\eusm Z$ dari
transformasi Fourier fungsi-fungsi dalam $\eusm D$. Keuntungan memulai dengan
$\eusm S$ adalah diperolehnya teori yang simetris, sebab
$\varhat{h}\in\eusm S$ untuk setiap $h\in\eusm S$. Namun mudah ditemukan objek
(misalnya fungsi $x\mapsto e^{x^2}$ atau $x\mapsto 1/|x|$) yang tidak dapat
ditafsirkan sebagai fungsional pada $\eusm S$, sehingga transformasi
Fouriernya, jika dapat dikaji, harus diselidiki dengan metode lain. Dalam 284Xq
saya menguraikan beberapa argumen yang membenarkan pernyataan bahwa transformasi
Fourier fungsi $x\mapsto 1/x$ adalah, atau dapat diwakili oleh, fungsi
$y\mapsto -i\sqrt{\bover{\pi}2}\sgn y$; asas umumnya di sini adalah mendekati
$0$ maupun $\infty$ secara simetris. Untuk berbagai pasangan serupa yang
ditetapkan melalui argumen berdasarkan gagasan 284Xr, lihat {\smc Lighthill
59}, bab.\ 3.

Karena itu tampaknya, setelah dua abad, % Fourier "Analytic Theory of Heat" 1822
%    1807 Mémoire sur la propagation de la chaleur dans les corps solides
% Wikipedia:  "When Fourier submitted a later competition essay in 1811,
% the committee (which included Lagrange, Laplace, Malus and Legendre,
% among others) concluded: ...the manner in which the author arrives at
% these equations is not exempt of difficulties and...his analysis to
% integrate them still leaves something to be desired on the score of
% generality and even rigour [citation needed]"
 kita masih harus melanjutkan dengan memeriksa secara cermat kelas-kelas fungsi
tertentu dan menguji interpretasi rumus yang sesuai. Dalam pembahasan di atas
saya berulang kali menggunakan konsep

\Centerline{$\lim_{a\to\infty}\int_{-a}^af$, \quad$\lim_{\epsilon\downarrow
0}\int_{-\infty}^{\infty} e^{-\epsilon x^2}f(x)dx$}

\noindent sebagai interpretasi alternatif bagi $\int_{-\infty}^{\infty}f$.
(Tentu saja keduanya berkaitan erat; lihat 284Xf.) Kernel khusus
$e^{-\epsilon x^2}$ digunakan karena termasuk $\eusm S$, merupakan fungsi
genap, transformasi Fouriernya dapat dihitung dan mudah dimanipulasi, serta
berkaitan dengan fungsi densitas probabilitas normal
$\bover1{\sigma\sqrt{2\pi}}e^{-x^2/2\sigma^2}$. Dengan demikian, fakta-fakta
yang kita peroleh mungkin berguna di tempat lain. Namun ada penerapan yang
lebih mudah ditangani dengan kernel lain -- misalnya,
$e^{-\epsilon|x|}$ (283Xq, 283Yc, 284Yf).)

Salah satu asas penuntun di sini ialah bahwa manipulasi yang murni formal,
seperti dalam daftar di atas, dan terutama perubahan urutan integrasi beserta
pertukaran limit lainnya, berulang kali menghasilkan rumus yang sah jika
ditafsirkan dengan tepat. Kuliah analisis pertama sering menanamkan hambatan:
mahasiswa diajari mencurigai setiap manipulasi yang belum dapat mereka
benarkan. Bagi saya, daya tarik pokok ini terutama terletak pada keragaman
argumen yang dituntut oleh pendekatan ketat, dengan landasan yang terus berubah
menurut konteks; tetapi optimisme yang ceria sering merupakan panduan terbaik
untuk memilih manipulasi yang patut dicoba untuk dibenarkan.

Karena buku ini membahas teori ukur, saya tentu sangat tertarik pada kemungkinan
suatu ukuran muncul sebagai transformasi Fourier. Inilah yang terjadi jika kita
mencari transformasi Fourier fungsi konstan $\chi\Bbb R$ (284R). Secara lebih
umum, setiap fungsi tempered periodik $f$ dengan periode $2\pi$ dapat diberi
transformasi Fourier berupa `ukuran bertanda' (untuk tujuan kita saat ini,
kombinasi linear kompleks dari ukuran-ukuran) yang terkonsentrasi pada
$\Bbb Z$. Massa pada setiap $k\in\Bbb Z$ ditentukan oleh koefisien Fourier yang
bersesuaian dari $f\restr\ocint{-\pi,\pi}$ (284Xo, 284Ya). Dalam bagian
berikutnya saya akan melangkah lebih jauh ke arah ini, terutama mengenai
distribusi probabilitas pada $\BbbR^r$. Alasan ukuran {\it positif} belum
mendesak untuk diperhatikan sejauh ini ialah bahwa kita tidak mengharapkan
fungsi positif sebagai transformasi Fourier kecuali jika beberapa syarat yang
sangat khusus dipenuhi, seperti dalam 283Yc.

Seperti dalam \S282, saya menggunakan struktur ruang Hilbert $L_{\Bbb C}^2$
sebagai dasar pembahasan transformasi Fourier fungsi-fungsi dalam
$\eusm L_{\Bbb C}^2$ (284O-284P). Namun, seperti untuk deret Fourier, teorema
Carleson (286U) memberikan uraian yang lebih langsung.

Dalam 284Wj, saya memberikan perhitungan berdasarkan rumus perubahan variabel
dalam 263A untuk menyajikan versi multidimensi dari pembalikan dan dilasi.
Namun tujuan sebenarnya ialah menunjukkan bahwa transformasi Fourier pada
$\BbbR^r$ didasarkan pada geometri hasil kali dalam Euklides, bukan pada sistem
koordinat Kartesius. }%end of notes

\discrpage
